\documentclass[12pt,a4paper]{article}
\usepackage[utf8]{inputenc}
\usepackage[T2A]{fontenc}
\usepackage[russian]{babel}
\usepackage{amsmath,amssymb}
\usepackage[left=25mm,right=25mm,top=20mm,bottom=20mm]{geometry}
\usepackage{enumitem}

\newcommand{\dd}{\mathrm{d}}
\newcommand{\pp}{\partial}
\newcommand{\ext}{\mathrm{ext}}
\newcommand{\eff}{\mathrm{eff}}
\newcommand{\st}{\text{ст}}
\newcommand{\sat}{\mathrm{sat}}

\title{Метод интегральных соотношений (моментов)\\
для задачи двухфазной фильтрации Black~Oil\\
с балансом энергии в радиальной геометрии}
\author{}
\date{}

\begin{document}
\maketitle
\tableofcontents
\newpage

%=============================================================================
\section{Введение и мотивация}\label{sec:intro}
%=============================================================================

Полная система уравнений двухфазной фильтрации Black~Oil с балансом
энергии выведена в~\cite{draft4}. Система состоит из трёх PDE
на $p(r,t)$, $T(r,t)$, $S_g(r,t)$ в области $r\in[r_w, R_\ext]$,
$t \geqslant 0$.

Стратегия приближённого решения: метод интегральных соотношений
(моментов) Баренблатта~\cite{barenblatt1954} применяется ко
\emph{всем трём уравнениям}. PDE сводятся к системе ODE для
параметров профилей $p(r,t)$, $S_g(r,t)$ и $T(r,t)$.
Температурное поле входит в уравнения давления и
насыщенности через PVT-свойства $R_s(p,T)$, $B_o(p,T)$,
$\mu_o(p,T)$ и т.\,д., а также через
$C_T\,\pp T/\pp t$ и $\beta_o^\sat\,\pp T/\pp t$,
поэтому <<последовательное>> решение
(сначала гидродинамика, потом температура) приводит
к неконтролируемым ошибкам.


%=============================================================================
\section{Исходные уравнения}\label{sec:equations}
%=============================================================================

Система состоит из трёх уравнений на три неизвестных
$p(r,t)$, $T(r,t)$, $S_g(r,t)$. Все PVT-свойства зависят
от $p$ и $T$: $R_s(p,T)$, $B_o(p,T)$, $B_g(p,T)$,
$\mu_\alpha(p,T)$ и т.\,д.

\subsection{Уравнение давления}

Из~\cite{draft4}, уравнение~(pressure):
\begin{equation}\label{eq:p_eq}
  \boxed{
    C_p\,\frac{\pp p}{\pp t}
    + C_T\,\frac{\pp T}{\pp t}
    + \mathcal{D} = 0,
  }
\end{equation}
где:

\paragraph{Подвижности фаз:}
\begin{equation}\label{eq:mobility}
  \lambda_\alpha = \frac{k\,k_{r\alpha}(S_g)}{\mu_\alpha(p,T)},
  \qquad \alpha \in \{o, g\}.
\end{equation}

\paragraph{Суммарная приведённая подвижность:}
\begin{equation}\label{eq:Lambda}
  \Lambda(p,T,S_g) = \frac{\lambda_o + \lambda_g}{B_g}
  = \frac{k}{B_g(p,T)}\!\left(
      \frac{k_{ro}(S_g)}{\mu_o(p,T)}
      + \frac{k_{rg}(S_g)}{\mu_g(p,T)}
    \right).
\end{equation}

\paragraph{Коэффициент ёмкости при $\pp p/\pp t$:}
\begin{equation}\label{eq:Cp}
  C_p = \phi\!\left[
    S_o\,(c_r + c_o^\sat)\!\left(\frac{1}{B_g} - \frac{R_s}{B_o}\right)
    + \xi_o\!\left(\frac{\pp R_s}{\pp p}
      + R_s\,(c_r + c_o^\sat)\right)
    + \frac{S_g}{B_g}\,(c_g + c_r)
  \right],
\end{equation}
где $\xi_o = S_o/B_o$, $S_o = 1 - S_g$, а <<полная>> сжимаемость
нефтяной фазы вдоль кривой насыщения:
\begin{equation}\label{eq:co_sat}
  c_o^\sat = c_o
    - \frac{1}{B_o}\frac{\pp B_o}{\pp R_s}\bigg|_{p,T}
      \frac{\pp R_s}{\pp p}\bigg|_T,
  \qquad
  c_o = -\frac{1}{B_o}\frac{\pp B_o}{\pp p}\bigg|_{T,\,R_s}.
\end{equation}

\paragraph{Коэффициент при $\pp T/\pp t$:}
\begin{equation}\label{eq:CT}
  C_T = \phi\!\left[
    -S_o\,\beta_o^\sat\!\left(\frac{1}{B_g} - \frac{R_s}{B_o}\right)
    + \xi_o\!\left(\frac{\pp R_s}{\pp T}
      - R_s\,\beta_o^\sat\right)
    - \frac{S_g}{B_g}\,\beta_g
  \right],
\end{equation}
где
\begin{equation}\label{eq:beta_o_sat}
  \beta_o^\sat = \beta_o
    + \frac{1}{B_o}\frac{\pp B_o}{\pp R_s}\bigg|_{p,T}
      \frac{\pp R_s}{\pp T}\bigg|_p,
  \qquad
  \beta_o = \frac{1}{B_o}\frac{\pp B_o}{\pp T}\bigg|_{p,\,R_s}.
\end{equation}

\paragraph{Пространственный оператор:}
\begin{equation}\label{eq:D_operator}
  \mathcal{D} = -\frac{1}{r}\,\frac{\pp}{\pp r}\!\left(
      r\,\Lambda\,\frac{\pp p}{\pp r}
    \right)
    - \frac{\lambda_o}{B_o}\!\left(
      \frac{\pp R_s}{\pp p}\,\frac{\pp p}{\pp r}
      + \frac{\pp R_s}{\pp T}\,\frac{\pp T}{\pp r}
    \right)\frac{\pp p}{\pp r}.
\end{equation}

\paragraph{Замечание.}
Знак при $\mathcal{D}$ в~\eqref{eq:p_eq}: $C_p\,\dot{p} + C_T\,\dot{T}
+ \mathcal{D} = 0$, т.\,е.\ $\mathcal{D}$ включает все
пространственные члены, перенесённые в левую часть.
В дивергентном слагаемом стоит минус, поскольку отток жидкости
($\mathcal{D} > 0$) снижает давление ($\dot{p} < 0$).


\subsection{Уравнение насыщенности}

Из~\cite{draft4}, уравнение~(Sg\_ode):
\begin{equation}\label{eq:Sg_eq}
  \boxed{
    \frac{\pp S_g}{\pp t}
    = S_o\,(c_r + c_o^\sat)\,\frac{\pp p}{\pp t}
      - S_o\,\beta_o^\sat\,\frac{\pp T}{\pp t}
      + \frac{B_o}{\phi}\,\frac{1}{r}\,\frac{\pp}{\pp r}\!
        \left(r\,\frac{u_o}{B_o}\right).
  }
\end{equation}
Подставляя $u_o = -\lambda_o\,\pp p/\pp r$:
\begin{equation}\label{eq:Sg_eq2}
  \frac{\pp S_g}{\pp t}
  = S_o\,(c_r + c_o^\sat)\,\frac{\pp p}{\pp t}
    - S_o\,\beta_o^\sat\,\frac{\pp T}{\pp t}
    - \frac{B_o}{\phi}\,\frac{1}{r}\,\frac{\pp}{\pp r}\!
      \left(\frac{r\,\lambda_o}{B_o}\,\frac{\pp p}{\pp r}\right).
\end{equation}

Член $-S_o\,\beta_o^\sat\,\pp T/\pp t$ описывает влияние
температуры на насыщенность: при охлаждении ($\pp T/\pp t < 0$)
растворимость газа растёт ($\pp R_s/\pp T < 0$ для типичных нефтей),
газ частично растворяется, и $S_g$ снижается. Знак зависит
от конкурирующих вкладов в $\beta_o^\sat$.


\subsection{Уравнение энергии}

Из~\cite{draft4}, уравнение~(energy\_final):
\begin{equation}\label{eq:energy}
  \boxed{
    (\rho c_p)_\eff\,\frac{\pp T}{\pp t}
    + (\rho c_p u)_\mathrm{conv}\,\frac{\pp T}{\pp r}
    = \frac{1}{r}\,\frac{\pp}{\pp r}\!\left(
        r\,\lambda_\eff\,\frac{\pp T}{\pp r}
      \right)
    + Q_\mathrm{JT} + Q_\mathrm{ad} + Q_\mathrm{deg},
  }
\end{equation}
где:

\paragraph{Эффективная теплоёмкость:}
\begin{equation}\label{eq:rhocp_eff}
  (\rho c_p)_\eff
  = (1-\phi)\,(\rho c_p)_s
    + \phi\,(S_o\,\rho_o\,c_{p,o} + S_g\,\rho_g\,c_{p,g}).
\end{equation}

\paragraph{Конвективный коэффициент:}
\begin{equation}\label{eq:rhocp_conv}
  (\rho c_p u)_\mathrm{conv}
  = \rho_o\,c_{p,o}\,u_o + \rho_g\,c_{p,g}\,u_g.
\end{equation}

\paragraph{Эффект Джоуля--Томсона:}
\begin{equation}\label{eq:QJT}
  Q_\mathrm{JT} = -\sum_\alpha u_\alpha\,(1 - \alpha_\alpha T)\,
    \frac{\pp p}{\pp r},
\end{equation}
где $\alpha_\alpha$ --- коэффициент теплового расширения
фазы~$\alpha$ при фиксированном составе.

\paragraph{Адиабатический эффект:}
\begin{equation}\label{eq:Qad}
  Q_\mathrm{ad} = \phi\!\left[
    T\sum_\alpha S_\alpha\,\alpha_\alpha
    + c_r\!\left(p + (\rho c_p)_s\,T\right)
  \right]\frac{\pp p}{\pp t}.
\end{equation}

\paragraph{Теплота дегазации:}
\begin{equation}\label{eq:Qdeg}
  Q_\mathrm{deg} = -L(p,T)\,\dot{m},
\end{equation}
где $L(p,T)$ --- удельная теплота дегазации, а массовый темп
выделения газа:
\begin{equation}\label{eq:mdot}
  \dot{m} = -\rho_g^\st\!\left[
    \phi\,\xi_o\!\left(
      \frac{\pp R_s}{\pp p}\,\frac{\pp p}{\pp t}
      + \frac{\pp R_s}{\pp T}\,\frac{\pp T}{\pp t}
    \right)
    + \frac{u_o}{B_o}\!\left(
      \frac{\pp R_s}{\pp p}\,\frac{\pp p}{\pp r}
      + \frac{\pp R_s}{\pp T}\,\frac{\pp T}{\pp r}
    \right)
  \right].
\end{equation}

\paragraph{Замечание.}
$Q_\mathrm{deg}$ содержит $\pp T/\pp t$ и $\pp T/\pp r$
через~\eqref{eq:mdot}. При подстановке в~\eqref{eq:energy}
они перегруппировываются с коэффициентами при $\pp T/\pp t$
и $\pp T/\pp r$:
\begin{equation}\label{eq:rhocp_eff_mod}
  (\rho c_p)_\eff^* = (\rho c_p)_\eff
    - L\,\rho_g^\st\,\phi\,\xi_o\,\frac{\pp R_s}{\pp T},
\end{equation}
\begin{equation}\label{eq:rhocp_conv_mod}
  (\rho c_p u)_\mathrm{conv}^* = (\rho c_p u)_\mathrm{conv}
    - L\,\rho_g^\st\,\frac{u_o}{B_o}\,\frac{\pp R_s}{\pp T}.
\end{equation}
Аналогично, часть $Q_\mathrm{deg}$, содержащая $\pp p/\pp t$
и $\pp p/\pp r$, поглощается в модифицированные
$Q_\mathrm{ad}^*$ и $Q_\mathrm{JT}^*$.
Перегруппированное уравнение энергии:
\begin{equation}\label{eq:energy_regrouped}
  (\rho c_p)_\eff^*\,\frac{\pp T}{\pp t}
  + (\rho c_p u)_\mathrm{conv}^*\,\frac{\pp T}{\pp r}
  = \frac{1}{r}\,\frac{\pp}{\pp r}\!\left(
      r\,\lambda_\eff\,\frac{\pp T}{\pp r}
    \right)
  + Q_\mathrm{JT}^* + Q_\mathrm{ad}^*,
\end{equation}
Явно:
\begin{equation}\label{eq:QJT_star}
  Q_\mathrm{JT}^*
  = Q_\mathrm{JT}
    + L\,\rho_g^\st\,\frac{u_o}{B_o}\,
      \frac{\pp R_s}{\pp p}\,\frac{\pp p}{\pp r},
\end{equation}
\begin{equation}\label{eq:Qad_star}
  Q_\mathrm{ad}^*
  = Q_\mathrm{ad}
    + L\,\rho_g^\st\,\phi\,\xi_o\,
      \frac{\pp R_s}{\pp p}\,\frac{\pp p}{\pp t}.
\end{equation}
Таким образом, $Q_\mathrm{JT}^*$ и $Q_\mathrm{ad}^*$ зависят
только от $p$, $T$, $S_g$ и производных давления (не от $\pp T$).


\subsection{Граничные и начальные условия}

\begin{equation}\label{eq:bc_p}
  p(r_w, t) = p_{wf}(t), \qquad
  p(R_\ext, t) = p_0.
\end{equation}
\begin{equation}\label{eq:bc_T}
  T(R_\ext, t) = T_0, \qquad
  \frac{\pp T}{\pp r}\bigg|_{r = r_w} = 0.
\end{equation}
\begin{equation}\label{eq:ic}
  p(r, 0) = p_0, \qquad T(r, 0) = T_0, \qquad S_g(r, 0) = 0.
\end{equation}

При $t = 0$ давление и температура однородны, газ полностью
растворён ($p_0 > p_b(T_0)$). При снижении $p_{wf} < p_b$
вблизи скважины возникает зона разгазирования
$r \in [r_w, r_b(t)]$, где $p < p_b(T)$.


%=============================================================================
\section{Определение интегральных моментов}\label{sec:moments_def}
%=============================================================================

\subsection{Зона разгазирования}

Радиус зоны разгазирования $r_b(t)$ определяется из:
\begin{equation}\label{eq:rb}
  p(r_b(t), t) = p_b(T(r_b(t), t)).
\end{equation}
При $r > r_b(t)$: $p > p_b(T)$, $S_g = 0$, фильтрация однофазная.

\subsection{Моменты давления}

\begin{equation}\label{eq:Mn}
  M_n(t) = \int_{r_w}^{R_\ext} r^{2n+1}\,p(r,t)\,\dd r,
  \qquad n = 0, 1, 2, \ldots
\end{equation}

\paragraph{Нулевой момент ($n = 0$):}
\begin{equation}\label{eq:M0}
  M_0(t) = \int_{r_w}^{R_\ext} r\,p(r,t)\,\dd r.
\end{equation}

\paragraph{Первый момент ($n = 1$):}
\begin{equation}\label{eq:M1}
  M_1(t) = \int_{r_w}^{R_\ext} r^3\,p(r,t)\,\dd r.
\end{equation}

\subsection{Момент насыщенности}

Газонасыщенность имеет компактный носитель $[r_w, r_b(t)]$:
\begin{equation}\label{eq:Sn}
  \mathcal{S}_n(t) = \int_{r_w}^{r_b(t)} r^{2n+1}\,S_g(r,t)\,\dd r.
\end{equation}

$\mathcal{S}_0$ пропорционален объёму свободного газа
в зоне разгазирования.

\subsection{Момент температуры}

Температурное поле определено на всей области $[r_w, R_\ext]$.
Нулевой момент:
\begin{equation}\label{eq:Tn}
  \mathcal{T}_0(t) = \int_{r_w}^{R_\ext} r\,T(r,t)\,\dd r.
\end{equation}

Физический смысл: $2\pi h\,\mathcal{T}_0$ пропорционален
полной тепловой энергии (отклонению от $T_0$) в пласте.


%=============================================================================
\section{Интегральные соотношения}\label{sec:integral_rel}
%=============================================================================

\subsection{Нулевой момент уравнения давления}

Умножим~\eqref{eq:p_eq} на $r$ и проинтегрируем по $r$
от $r_w$ до $R_\ext$. Используя $\mathcal{D}$
из~\eqref{eq:D_operator}:
\begin{multline}\label{eq:moment0}
  \boxed{
    \int_{r_w}^{R_\ext} r\,C_p\,\frac{\pp p}{\pp t}\,\dd r
    + \int_{r_w}^{R_\ext} r\,C_T\,\frac{\pp T}{\pp t}\,\dd r
  } \\
  \boxed{
    = \left[r\,\Lambda\,\frac{\pp p}{\pp r}\right]_{r_w}^{R_\ext}
    + \int_{r_w}^{R_\ext}
      r\,\frac{\lambda_o}{B_o}\!\left(
        \frac{\pp R_s}{\pp p}\,\frac{\pp p}{\pp r}
        + \frac{\pp R_s}{\pp T}\,\frac{\pp T}{\pp r}
      \right)\frac{\pp p}{\pp r}\,\dd r.
  }
\end{multline}

По сравнению с изотермическим случаем появились:
\begin{itemize}[nosep]
  \item член $\int r\,C_T\,(\pp T/\pp t)\,\dd r$ в левой части;
  \item конвективный перенос газа за счёт градиента температуры
        $(\pp R_s/\pp T)\,(\pp T/\pp r)\,(\pp p/\pp r)$ в квадратичном
        интеграле.
\end{itemize}

\paragraph{Вычисление квадратичного интеграла.}
В зоне разгазирования координаты $\eta$ и $\zeta$
связаны линейно: $\zeta = 1 - (L_b/L_0)\,\eta$.
Поэтому произведение $(\pp p/\pp r)\,(\pp T/\pp r)$
при замороженных PVT-свойствах выражается через степенные
функции $\eta$:
\[
  \frac{\pp T}{\pp r}\,\frac{\pp p}{\pp r}
  = \frac{2\sigma\,\Delta T_w\,\Delta p_1\,\gamma}
    {r^2\,L_0\,L_b}\,
    \zeta^{\sigma-1}(1-\zeta^\sigma)\,\eta^{\gamma-1},
\]
где $\zeta(\eta) = 1-(L_b/L_0)\eta$.
При $L_b \ll L_0$ (типично: $L_b/L_0 \sim 0{,}01$--$0{,}1$)
$\zeta \approx 1$ в зоне разгазирования, и
интеграл упрощается.
В однофазной зоне ($\eta$ не определена, но $S_g = 0$)
вклад $\pp R_s/\pp p = 0$ при $p > p_b^*$, и
квадратичный интеграл обнуляется.

\subsection{Первый момент уравнения давления}

Умножим~\eqref{eq:p_eq} на $r^3$ и проинтегрируем:
\begin{multline}\label{eq:moment1}
  \int_{r_w}^{R_\ext} r^3\,C_p\,\frac{\pp p}{\pp t}\,\dd r
  + \int_{r_w}^{R_\ext} r^3\,C_T\,\frac{\pp T}{\pp t}\,\dd r \\
  = \left[r^3\,\Lambda\,\frac{\pp p}{\pp r}\right]_{r_w}^{R_\ext}
    - 2\int_{r_w}^{R_\ext} r^2\,\Lambda\,\frac{\pp p}{\pp r}\,\dd r
  + \int_{r_w}^{R_\ext}
    r^3\,\frac{\lambda_o}{B_o}\!\left(
      \frac{\pp R_s}{\pp p}\,\frac{\pp p}{\pp r}
      + \frac{\pp R_s}{\pp T}\,\frac{\pp T}{\pp r}
    \right)\frac{\pp p}{\pp r}\,\dd r.
\end{multline}

\subsection{Нулевой момент уравнения насыщенности}

Умножим~\eqref{eq:Sg_eq2} на $r$ и проинтегрируем от $r_w$ до $r_b$:
\begin{multline}\label{eq:Sg_moment0}
  \frac{\dd\mathcal{S}_0}{\dd t}
  = \int_{r_w}^{r_b} r\,S_o\,(c_r+c_o^\sat)\,
      \frac{\pp p}{\pp t}\,\dd r
    - \int_{r_w}^{r_b} r\,S_o\,\beta_o^\sat\,
      \frac{\pp T}{\pp t}\,\dd r \\
  - \frac{B_o}{\phi}\!\left[
      \frac{r\,\lambda_o}{B_o}\,\frac{\pp p}{\pp r}
    \right]_{r_w}^{r_b}.
\end{multline}

Член с $\dot{r}_b$ обращается в нуль ($S_g(r_b,t) = 0$,
правило Лейбница). Граничные потоки:
\begin{equation}\label{eq:Sg_flux}
  \left[\frac{r\,\lambda_o}{B_o}\,\frac{\pp p}{\pp r}\right]_{r_w}^{r_b}
  = \frac{r_b\,\lambda_o(r_b)}{B_o(p_b,T_b)}\,p'_b
    - \frac{r_w\,\lambda_o(r_w)}{B_o(p_{wf},T_w)}\,p'_w,
\end{equation}
где $T_b = T(r_b,t)$, $T_w = T(r_w,t)$.

По сравнению с изотермическим случаем появился
интеграл $-\int r\,S_o\,\beta_o^\sat\,(\pp T/\pp t)\,\dd r$.


\subsection{Нулевой момент уравнения энергии}

Умножим~\eqref{eq:energy_regrouped} на $r$ и проинтегрируем
от $r_w$ до $R_\ext$:
\begin{multline}\label{eq:T_moment0}
  \int_{r_w}^{R_\ext} r\,(\rho c_p)_\eff^*\,
    \frac{\pp T}{\pp t}\,\dd r
  + \int_{r_w}^{R_\ext} r\,(\rho c_p u)_\mathrm{conv}^*\,
    \frac{\pp T}{\pp r}\,\dd r \\
  = \left[r\,\lambda_\eff\,\frac{\pp T}{\pp r}
    \right]_{r_w}^{R_\ext}
  + \int_{r_w}^{R_\ext} r\,(Q_\mathrm{JT}^* + Q_\mathrm{ad}^*)\,\dd r.
\end{multline}

Дивергентная часть вычисляется точно:
\begin{equation}\label{eq:T_flux}
  \left[r\,\lambda_\eff\,\frac{\pp T}{\pp r}
    \right]_{r_w}^{R_\ext}
  = R_\ext\,\lambda_\eff\,T'_\ext
    - \underbrace{r_w\,\lambda_\eff\,T'_w}_{= 0},
\end{equation}
поскольку граничное условие~\eqref{eq:bc_T} даёт $T'_w = 0$.


%=============================================================================
\section{Параметризация профилей}\label{sec:profiles}
%=============================================================================

\subsection{Структура области}

Область $[r_w, R_\ext]$ делится на две зоны:
\begin{itemize}[nosep]
  \item \textbf{Зона разгазирования} $r \in [r_w, r_b(t)]$:
        $p < p_b(T)$, $S_g > 0$, двухфазная фильтрация.
  \item \textbf{Однофазная зона} $r \in [r_b(t), R_\ext]$:
        $p \geqslant p_b(T)$, $S_g = 0$, однофазная фильтрация нефти.
\end{itemize}

\subsection{Профиль давления}

\paragraph{Однофазная зона ($r_b \leqslant r \leqslant R_\ext$).}
Квазистационарный логарифмический профиль:
\begin{equation}\label{eq:p_outer}
  p(r,t) = p_b^* + (p_0 - p_b^*)\,
    \frac{\ln(r/r_b)}{\ln(R_\ext/r_b)},
  \qquad r_b \leqslant r \leqslant R_\ext.
\end{equation}

\paragraph{Зона разгазирования ($r_w \leqslant r \leqslant r_b$).}
Степенной профиль:
\begin{equation}\label{eq:p_inner}
  p(r,t) = p_{wf} + (p_b^* - p_{wf})\,\Phi\!\left(
    \frac{\ln(r/r_w)}{\ln(r_b/r_w)}
  \right),
\end{equation}
\begin{equation}\label{eq:Phi}
  \Phi(\eta) = \eta^\gamma, \qquad
  \eta = \frac{\ln(r/r_w)}{\ln(r_b/r_w)},
  \qquad \gamma > 0.
\end{equation}

\paragraph{Непрерывность потока на $r = r_b$.}
При $S_g(r_b) = 0$ подвижности совпадают.
Обозначим $p_b^* = p_b(T(r_b))$ --- давление насыщения
при температуре на фронте разгазирования.
Тогда $\Delta p_1 = p_b^* - p_{wf}$,
$\Delta p_2 = p_0 - p_b^*$, и
\begin{equation}\label{eq:gamma_constraint}
  \gamma = \frac{\Delta p_2\,L_b}{\Delta p_1\,L_e}.
\end{equation}
Величина $p_b^* = p_b(T(r_b))$ зависит от $\Delta T_w$
через профиль~\eqref{eq:T_profile}:
\begin{equation}\label{eq:pb_star}
  p_b^* = p_b\!\left(T_0 - \Delta T_w\,\Psi(\zeta_b)\right),
  \qquad
  \zeta_b = \frac{\ln(R_\ext/r_b)}{L_0} = \frac{L_e}{L_0}.
\end{equation}
При $\Delta T_w = 0$ (начальная стадия):
$p_b^* = p_b(T_0)$ --- фиксированное значение.

Производные давления:
\begin{equation}\label{eq:dp_inner}
  \frac{\pp p}{\pp r}\bigg|_\text{inner}
  = \frac{\Delta p_1\,\gamma}{r\,L_b}\,\eta^{\gamma-1},
\end{equation}
\begin{equation}\label{eq:dp_outer}
  \frac{\pp p}{\pp r}\bigg|_\text{outer}
  = \frac{\Delta p_2}{r\,L_e}.
\end{equation}

\paragraph{Обозначения.}
$L_b = \ln(r_b/r_w)$, $L_e = \ln(R_\ext/r_b)$,
$\Delta p_1 = p_b^* - p_{wf}$, $\Delta p_2 = p_0 - p_b^*$
(здесь и далее $p_b^* = p_b(T(r_b))$, см.~\eqref{eq:pb_star}).


\subsection{Профиль насыщенности}

Газонасыщенность имеет компактный носитель $[r_w, r_b(t)]$:
\begin{equation}\label{eq:Sg_profile}
  S_g(r,t) = S_{g,w}(t)\!\left[
    1 - \left(\frac{\ln(r/r_w)}{\ln(r_b/r_w)}\right)^{\!\beta}
  \right],
  \qquad r_w \leqslant r \leqslant r_b,
\end{equation}
$S_g(r_w) = S_{g,w}$, $S_g(r_b) = 0$, $\beta > 0$.


\subsection{Профиль температуры}\label{ssec:T_profile}

Температурное поле определено на $[r_w, R_\ext]$ и удовлетворяет:
\begin{itemize}[nosep]
  \item $T(R_\ext, t) = T_0$ (контур питания);
  \item $T'(r_w, t) = 0$ (адиабатическая стенка).
\end{itemize}

Профиль параметризуем функцией формы
с амплитудным параметром $\Delta T_w(t) = T_0 - T_w(t)$,
где $T_w(t) = T(r_w, t)$ --- температура на забое:
\begin{equation}\label{eq:T_profile}
  T(r,t) = T_0 - \Delta T_w(t)\,\Psi(\zeta),
  \qquad
  \zeta = \frac{\ln(R_\ext/r)}{\ln(R_\ext/r_w)} = \frac{\ln(R_\ext/r)}{L_0}.
\end{equation}

Функция формы $\Psi(\zeta)$ должна удовлетворять:
$\Psi(0) = 0$ (контур питания),
$\Psi(1) = 1$ (забой),
$\Psi'(1) = 0$ (адиабатическая стенка: $\pp T/\pp r|_{r_w} = 0$).

Простейший профиль с этими свойствами:
\begin{equation}\label{eq:Psi}
  \boxed{
    \Psi(\zeta) = \zeta^\sigma(2 - \zeta^\sigma),
    \qquad \sigma > 0.
  }
\end{equation}

Проверка:
\begin{itemize}[nosep]
  \item $\Psi(0) = 0$. \checkmark
  \item $\Psi(1) = 1\cdot(2-1) = 1$. \checkmark
  \item $\Psi'(\zeta) = 2\sigma\zeta^{\sigma-1}(1 - \zeta^\sigma)$,
        поэтому $\Psi'(1) = 2\sigma\cdot 1\cdot 0 = 0$. \checkmark
\end{itemize}

Производная температуры:
\begin{equation}\label{eq:dT_dr}
  \frac{\pp T}{\pp r}
  = \frac{\Delta T_w}{r\,L_0}\,\Psi'(\zeta)
  = \frac{2\sigma\,\Delta T_w}{r\,L_0}\,
    \zeta^{\sigma-1}(1-\zeta^\sigma).
\end{equation}
При $r = r_w$ ($\zeta = 1$): $\pp T/\pp r = 0$. \checkmark
Максимум $|\pp T/\pp r|$ достигается внутри области.

\paragraph{Интегралы.}
Моменты профиля выражаются через:
\begin{equation}\label{eq:IPsi}
  J^{(m)} = \int_0^1 \zeta^m\,(2 - \zeta^\sigma)\,\dd\zeta
  = \frac{2}{m+1} - \frac{1}{m+\sigma+1}.
\end{equation}

\paragraph{Выбор $\sigma$.}
При преобладании конвекции ($\mathrm{Pe} \gg 1$)
нефтяной поток подводит невозмущённую
температуру $T_0$ от контура, а охлаждение
генерируется локально (JT-эффект, дегазация).
Возмущение $\Delta T$ концентрируется
вблизи $r_w$. При $\sigma = 1$: $\Psi = \zeta(2-\zeta) = 1-(1-\zeta)^2$ ---
параболический профиль, основное изменение
вблизи забоя. При $\sigma > 1$ профиль
ещё сильнее локализован у $r_w$.

$\sigma$ можно зафиксировать ($\sigma = 1$) или определять
из первого момента уравнения энергии (добавив четвёртое ODE).
Начинаем с фиксированного $\sigma = 1$.


%=============================================================================
\section{Система ODE}\label{sec:ode}
%=============================================================================

\subsection{Свободные параметры}

Профили содержат зависящие от времени параметры:
\begin{center}
\begin{tabular}{cl}
  $r_b(t)$       & радиус зоны разгазирования, \\
  $\gamma(t)$    & показатель формы давления
                    ($\gamma = \gamma(r_b)$
                    из~\eqref{eq:gamma_constraint}), \\
  $S_{g,w}(t)$   & газонасыщенность на забое, \\
  $\Delta T_w(t)$ & отклонение температуры на забое от $T_0$.
\end{tabular}
\end{center}

Показатели формы $\beta$ и $\sigma$ фиксированы.
$\gamma$ --- не независимый: $\gamma = \gamma(r_b)$.
Остаются \textbf{три свободных параметра}:
$r_b(t)$, $S_{g,w}(t)$, $\Delta T_w(t)$.

\subsection{Уравнения для параметров}

Подставляя профили~\eqref{eq:p_outer}--\eqref{eq:T_profile}
в три интегральных соотношения~\eqref{eq:moment0},
\eqref{eq:Sg_moment0}, \eqref{eq:T_moment0},
получаем систему трёх ODE:

\paragraph{1. Из нулевого момента давления~\eqref{eq:moment0}:}
\begin{equation}\label{eq:ode1}
  \mathcal{A}_{11}\,\dot{r}_b
  + \mathcal{A}_{12}\,\dot{S}_{g,w}
  + \mathcal{A}_{13}\,\dot{\Delta T}_w
  = \mathcal{F}_1.
\end{equation}

\paragraph{2. Из нулевого момента насыщенности~\eqref{eq:Sg_moment0}:}
\begin{equation}\label{eq:ode2}
  \mathcal{A}_{21}\,\dot{r}_b
  + \mathcal{A}_{22}\,\dot{S}_{g,w}
  + \mathcal{A}_{23}\,\dot{\Delta T}_w
  = \mathcal{F}_2.
\end{equation}

\paragraph{3. Из нулевого момента энергии~\eqref{eq:T_moment0}:}
\begin{equation}\label{eq:ode3}
  \mathcal{A}_{31}\,\dot{r}_b
  + \mathcal{A}_{32}\,\dot{S}_{g,w}
  + \mathcal{A}_{33}\,\dot{\Delta T}_w
  = \mathcal{F}_3.
\end{equation}

Матричная форма:
\begin{equation}\label{eq:ode_system}
  \boxed{
    \begin{pmatrix}
      \mathcal{A}_{11} & \mathcal{A}_{12} & \mathcal{A}_{13} \\
      \mathcal{A}_{21} & \mathcal{A}_{22} & \mathcal{A}_{23} \\
      \mathcal{A}_{31} & \mathcal{A}_{32} & \mathcal{A}_{33}
    \end{pmatrix}
    \begin{pmatrix}
      \dot{r}_b \\ \dot{S}_{g,w} \\ \dot{\Delta T}_w
    \end{pmatrix}
    =
    \begin{pmatrix}
      \mathcal{F}_1 \\ \mathcal{F}_2 \\ \mathcal{F}_3
    \end{pmatrix}.
  }
\end{equation}

\subsection{Структура коэффициентов}

Коэффициенты $\mathcal{A}_{ij}$ и $\mathcal{F}_i$ вычисляются
подстановкой профилей~\eqref{eq:p_outer}--\eqref{eq:T_profile}
в интегралы и дифференцированием по параметрам.

\paragraph{Связь параметров с производными профилей.}
Профили $p(r,t)$ и $T(r,t)$ зависят от $t$ через
$r_b(t)$, $\Delta T_w(t)$ (и $p_{wf}(t)$, задаваемое).
Поэтому $\pp p/\pp t$ и $\pp T/\pp t$ выражаются
через $\dot{r}_b$ и $\dot{\Delta T}_w$:
\begin{equation}\label{eq:dp_dt}
  \frac{\pp p}{\pp t}
  \approx \frac{\pp p}{\pp r_b}\,\dot{r}_b
    + \frac{\pp p}{\pp p_{wf}}\,\dot{p}_{wf},
\end{equation}
где опущен член
$(\pp p/\pp\Delta T_w)\,\dot{\Delta T}_w$,
пренебрежимо малый по сравнению с температурным
вкладом $C_T$ (обоснование в~\S\ref{sssec:profile_derivatives}).
\begin{equation}\label{eq:dT_dt}
  \frac{\pp T}{\pp t}
  = \frac{\pp T}{\pp r_b}\,\dot{r}_b
    + \frac{\pp T}{\pp \Delta T_w}\,\dot{\Delta T}_w.
\end{equation}
Профиль $S_g$ зависит от $r_b$ и $S_{g,w}$, поэтому
$\pp S_g/\pp t$ содержит $\dot{r}_b$ и $\dot{S}_{g,w}$.

Подставляя~\eqref{eq:dp_dt}--\eqref{eq:dT_dt}
в~\eqref{eq:moment0}--\eqref{eq:T_moment0}, собираем
коэффициенты при $\dot{r}_b$, $\dot{S}_{g,w}$,
$\dot{\Delta T}_w$ --- это элементы матрицы
$\mathcal{A}_{ij}$. Члены, не содержащие производных
по времени, образуют правую часть $\mathcal{F}_i$.

\paragraph{Температурный вклад в уравнение давления.}
Элемент $\mathcal{A}_{13}$ содержит:
\[
  \mathcal{A}_{13} \sim \int_{r_w}^{R_\ext} r\,C_T\,
    \frac{\pp T}{\pp \Delta T_w}\,\dd r.
\]
Из~\eqref{eq:T_profile}: $\pp T/\pp \Delta T_w = -\Psi(\zeta)$.
При $C_T < 0$ (типичный случай: $\pp R_s/\pp T < 0$)
и $\Psi > 0$: вклад положителен ($\mathcal{A}_{13} > 0$).
Влияние на динамику $r_b$ определяется из полной
системы~\eqref{eq:ode_system}.

\paragraph{Температурный вклад в уравнение насыщенности.}
Элемент $\mathcal{A}_{23}$ содержит (подробный
вывод~--- в~\S\ref{sssec:gas_mass_coefficients}):
\[
  \mathcal{A}_{23} \sim +\int_{r_w}^{r_b} r\,S_o\,\beta_o^\sat\,
    \frac{\pp T}{\pp \Delta T_w}\,\dd r
  = -\int_{r_w}^{r_b} r\,S_o\,\beta_o^\sat\,\Psi(\zeta)\,\dd r.
\]


\subsection{Явный вид коэффициентов}
\label{ssec:explicit_coefficients}

\paragraph{Двухзонное замораживание PVT-свойств.}
\label{par:frozen_pvt}
Для вычисления интегралов PVT-свойства замораживаются
раздельно в каждой зоне:
\begin{itemize}[nosep]
  \item \textbf{Зона разгазирования} ($r_w \leqslant r \leqslant r_b$):
        \[
          \bar{p}_1 = \frac{p_{wf} + p_b^*}{2}, \qquad
          \bar{T}_1 = T_0 - \Delta T_w\,\bar{\Psi}_1,
          \qquad
          \bar{S}_{g,1} = \frac{S_{g,w}}{2},
        \]
        где
        \[
          \bar{\Psi}_1
          = \frac{1}{L_b}\int_{\zeta_b}^{1}\Psi(\zeta)\,\dd\zeta.
        \]
        При $L_b \ll L_0$ ($\zeta_b \to 1$):
        $\bar{\Psi}_1 \to 1$, т.\,е.\ $\bar{T}_1 \approx T_0 - \Delta T_w$
        (вся зона разгазирования лежит вблизи забоя).
  \item \textbf{Однофазная зона} ($r_b \leqslant r \leqslant R_\ext$):
        \[
          \bar{p}_2 = \frac{p_b^* + p_0}{2}, \qquad
          \bar{T}_2 \approx T_0, \qquad
          S_g = 0.
        \]
\end{itemize}

Для компактности введём обозначения:
\begin{align}
  I_\Phi^{(n)} &= \int_0^1 \eta^{n\gamma}\,\dd\eta
    = \frac{1}{n\gamma + 1}, \label{eq:IPhi} \\
  I_S^{(n)} &= \int_0^1 (1-\eta^\beta)\,\eta^{n\gamma}\,\dd\eta
    = I_\Phi^{(n)} - \frac{1}{n\gamma + \beta + 1}.
    \label{eq:IS}
\end{align}
Интегралы профиля температуры $J^{(m)}$
определены в~\eqref{eq:IPsi}.

\paragraph{Расщепление матрицы.}
Профили $p(r)$ и $T(r)$ не зависят от $S_{g,w}$
(давление определяется через $r_b$, $p_{wf}$;
температура --- через $r_b$, $\Delta T_w$).
Поэтому $\pp p/\pp t$ и $\pp T/\pp t$ не содержат
$\dot{S}_{g,w}$, и
\begin{equation}\label{eq:matrix_structure}
  \mathcal{A}_{12} = 0, \qquad \mathcal{A}_{32} = 0.
\end{equation}
Матрица системы~\eqref{eq:ode_system} имеет вид:
\[
  \begin{pmatrix}
    \mathcal{A}_{11} & 0 & \mathcal{A}_{13} \\
    \mathcal{A}_{21} & \mathcal{A}_{22} & \mathcal{A}_{23} \\
    \mathcal{A}_{31} & 0 & \mathcal{A}_{33}
  \end{pmatrix}.
\]
Из первого и третьего уравнений определяются $\dot{r}_b$
и $\dot{\Delta T}_w$ (подсистема $2\times 2$), после чего
из второго уравнения находится $\dot{S}_{g,w}$.

\paragraph{Замечание.}
Расщепление касается только производных $\dot{S}_{g,w}$.
Величина $S_{g,w}$ входит в коэффициенты $\mathcal{A}_{ij}$
и $\mathcal{F}_i$ через PVT-свойства (оцениваемые
при замороженных значениях $\bar{S}_g$), поэтому
система остаётся связанной.

Ниже приводятся явные формулы для первой строки матрицы
($\mathcal{A}_{11}$, $\mathcal{A}_{13}$, $\mathcal{F}_1$)
--- уравнение сохранения массы нефти.%
\footnote{FEAT-005a: коэффициенты уравнения массы нефти.}

\subsubsection{Производные профилей по параметрам}
\label{sssec:profile_derivatives}

Для вычисления $\mathcal{A}_{ij}$ необходимы частные
производные профилей давления и температуры по
свободным параметрам $r_b$, $\Delta T_w$, $p_{wf}$.

\paragraph{Температурный профиль.}
Координата $\zeta = \ln(R_\ext/r)/L_0$ не зависит от $r_b$,
поэтому
\begin{equation}\label{eq:dT_drb}
  \frac{\pp T}{\pp r_b} = 0.
\end{equation}
Производная по амплитуде:
\begin{equation}\label{eq:dT_dTw}
  \frac{\pp T}{\pp \Delta T_w} = -\Psi(\zeta).
\end{equation}

\paragraph{Давление в зоне разгазирования.}
Из~\eqref{eq:p_inner}: $p = p_{wf} + \Delta p_1\,\eta^\gamma$.
При фиксированном $r$ и изменении $r_b$ зависят:
$\eta = \ln(r/r_w)/L_b$, $L_b = \ln(r_b/r_w)$,
$\gamma$~\eqref{eq:gamma_constraint},
$\Delta p_1 = p_b^* - p_{wf}$.
Частные производные:
\begin{equation}\label{eq:deta_drb}
  \frac{\pp\eta}{\pp r_b} = -\frac{\eta}{r_b\,L_b},
  \qquad
  \frac{\pp L_b}{\pp r_b} = \frac{1}{r_b}.
\end{equation}
Давление насыщения на фронте зависит от $r_b$
через температуру:
\begin{equation}\label{eq:dpb_drb}
  \frac{\pp p_b^*}{\pp r_b}
  = \frac{\dd p_b}{\dd T}\bigg|_{T(r_b)}
    \cdot \frac{\Delta T_w\,\Psi'(\zeta_b)}{r_b\,L_0} \geqslant 0.
\end{equation}
Знак: рост $r_b$ уменьшает $\zeta_b$ (фронт сдвигается
в более тёплую область) $\Rightarrow$ $T(r_b)$ растёт
$\Rightarrow$ $p_b^*$ растёт ($\dd p_b/\dd T > 0$
по Стендингу). При $\Delta T_w = 0$: $\pp p_b^*/\pp r_b = 0$
(температурное поле однородно, $p_b^*$ не зависит от $r_b$).

Полная производная давления в зоне~1:
\begin{equation}\label{eq:dp_drb_inner}
  \frac{\pp p}{\pp r_b}\bigg|_\text{inner}
  = \frac{\pp\Delta p_1}{\pp r_b}\,\eta^\gamma
    + \Delta p_1\,\eta^\gamma\!\left(
      \ln\eta\cdot\frac{\pp\gamma}{\pp r_b}
      + \frac{\gamma}{\eta}\,\frac{\pp\eta}{\pp r_b}
    \right),
\end{equation}
где $\pp\Delta p_1/\pp r_b = \pp p_b^*/\pp r_b$~\eqref{eq:dpb_drb}.
Производная $\pp\gamma/\pp r_b$ определяется
дифференцированием~\eqref{eq:gamma_constraint}
и зависит от $\pp p_b^*/\pp r_b$, $\pp L_b/\pp r_b$,
$\pp L_e/\pp r_b$. Формула~\eqref{eq:dp_drb_inner}
приводится для анализа структуры; при численной
реализации $\pp p/\pp r_b$ вычисляется конечной
разностью по $r_b$ (с пересчётом $p_b^*$ и $\gamma$
при каждом $r_b \pm \delta r_b$).

\paragraph{Давление в однофазной зоне.}
Из~\eqref{eq:p_outer}:
$p = p_b^* + \Delta p_2\,\ln(r/r_b)/L_e$,
где $\Delta p_2 = p_0 - p_b^*$, $L_e = \ln(R_\ext/r_b)$.
\begin{multline}\label{eq:dp_drb_outer}
  \frac{\pp p}{\pp r_b}\bigg|_\text{outer}
  = \frac{\pp p_b^*}{\pp r_b}
    \left(1 - \frac{\ln(r/r_b)}{L_e}\right) \\
    + \frac{\Delta p_2}{r_b}
    \left(
      \frac{\ln(r/r_b)}{L_e^2} - \frac{1}{L_e}
    \right).
\end{multline}
Предельные значения:
при $r = r_b$: $\pp p/\pp r_b = \pp p_b^*/\pp r_b
  - \Delta p_2/(r_b L_e)$;
при $r = R_\ext$: $\pp p/\pp r_b = 0$
(давление на контуре фиксировано $p = p_0$).

\paragraph{Производная по $p_{wf}$.}
В зоне разгазирования:
\begin{equation}\label{eq:dp_dpwf}
  \frac{\pp p}{\pp p_{wf}}\bigg|_\text{inner}
  = 1 - \eta^\gamma + \gamma\,\eta^\gamma\ln\eta,
\end{equation}
где учтено $\pp\gamma/\pp p_{wf} = \gamma/\Delta p_1$
из~\eqref{eq:gamma_constraint}.
Пределы: при $\eta = 0$ --- $\pp p/\pp p_{wf} = 1$
(давление на забое равно $p_{wf}$);
при $\eta = 1$ --- $\pp p/\pp p_{wf} = 0$.
В однофазной зоне: $\pp p/\pp p_{wf} = 0$
(давление $p_0$ и $p_b^*$ не зависят от $p_{wf}$).

\paragraph{Замечание о $\pp p/\pp\Delta T_w$.}
Давление формально зависит от $\Delta T_w$ через
$p_b^*(T(r_b))$~\eqref{eq:pb_star}. Однако вклад
$C_p\cdot\pp p/\pp\Delta T_w$ мал по сравнению с $C_T$.
В однофазной зоне ($S_g = 0$, $p > p_b$):
$C_p \approx \phi\,(c_r + c_o)/B_g \sim 10^{-8}$ [1/Па] (при $c_o = 0$ по Стендингу),
$\pp p/\pp\Delta T_w \sim (\dd p_b/\dd T)\,\Psi(\zeta_b)
  \sim 5\times 10^{4}$ [Па/К],
$C_T \approx -\phi\,\beta_o/B_g \sim -2\times 10^{-2}$ [1/К]:
\[
  \frac{C_p\cdot|\pp p/\pp\Delta T_w|}{|C_T|}
  \sim \frac{10^{-8}\cdot 5\times 10^4}{2\times 10^{-2}}
  \sim 2{,}5\times 10^{-2} \approx 2{,}5\%.
\]
Пренебрегается в соответствии с~\eqref{eq:dp_dt}.

\subsubsection{Коэффициенты первого уравнения (масса нефти)}
\label{sssec:oil_mass_coefficients}

Подставляя~\eqref{eq:dp_dt}--\eqref{eq:dT_dt}
в~\eqref{eq:moment0}, собираем коэффициенты
при $\dot{r}_b$, $\dot{S}_{g,w}$, $\dot{\Delta T}_w$.

\paragraph{Элемент $\mathcal{A}_{11}$.}
Коэффициент при $\dot{r}_b$:
\begin{equation}\label{eq:A11}
  \boxed{
    \mathcal{A}_{11}
    = \int_{r_w}^{R_\ext} r\,C_p\,
      \frac{\pp p}{\pp r_b}\,\dd r.
  }
\end{equation}
Член с $C_T$ обнуляется, т.\,к.\ $\pp T/\pp r_b = 0$~\eqref{eq:dT_drb}.

\emph{Знак:} $C_p > 0$ (ёмкость среды),
$\pp p/\pp r_b < 0$ (расширение зоны разгазирования
распределяет перепад давления на большую область,
снижая $p$ в каждой точке).
Поэтому $\mathcal{A}_{11} < 0$.

Интеграл разбивается на две зоны с замороженными PVT
($\bar{C}_p^{(1)}$, $\bar{C}_p^{(2)}$):
\begin{equation}\label{eq:A11_zones}
  \mathcal{A}_{11}
  = \bar{C}_p^{(1)}\!\int_{r_w}^{r_b} r\,
      \frac{\pp p}{\pp r_b}\bigg|_\text{inner}\dd r
    + \bar{C}_p^{(2)}\!\int_{r_b}^{R_\ext} r\,
      \frac{\pp p}{\pp r_b}\bigg|_\text{outer}\dd r,
\end{equation}
где $\pp p/\pp r_b$ в каждой зоне
даётся~\eqref{eq:dp_drb_inner}
и~\eqref{eq:dp_drb_outer}.
Экспоненциальные весовые функции $r\,\dd r
= r_w^2 L_b\,e^{2\eta L_b}\,\dd\eta$ не сводятся к
$I_\Phi^{(n)}$~\eqref{eq:IPhi}; интегралы вычисляются
численно.

\paragraph{Элемент $\mathcal{A}_{12} = 0$.}
Следует из расщепления~\eqref{eq:matrix_structure}:
давление не зависит от $S_{g,w}$.

\paragraph{Элемент $\mathcal{A}_{13}$.}
Коэффициент при $\dot{\Delta T}_w$:
\begin{equation}\label{eq:A13}
  \boxed{
    \mathcal{A}_{13}
    = \int_{r_w}^{R_\ext} r\,C_T\,
      \frac{\pp T}{\pp\Delta T_w}\,\dd r
    = -\int_{r_w}^{R_\ext} r\,C_T\,\Psi(\zeta)\,\dd r.
  }
\end{equation}
Двухзонное разбиение:
\begin{equation}\label{eq:A13_zones}
  \mathcal{A}_{13}
  = -\bar{C}_T^{(1)}\!\int_{r_w}^{r_b} r\,\Psi(\zeta)\,\dd r
    - \bar{C}_T^{(2)}\!\int_{r_b}^{R_\ext} r\,\Psi(\zeta)\,\dd r.
\end{equation}
Интегралы содержат $r\,\dd r = R_\ext^2 L_0\,e^{-2\zeta L_0}\,\dd\zeta$
и не сводятся к $J^{(m)}$~\eqref{eq:IPsi}; также вычисляются численно.

\emph{Знак:} при $C_T < 0$ и $\Psi > 0$: $\mathcal{A}_{13} > 0$.
Математически: $-C_T > 0$ и $\Psi > 0$, поэтому
подынтегральная в~\eqref{eq:A13} положительна.
Влияние на динамику $r_b$ определяется
решением полной системы~\eqref{eq:ode_system}.

\paragraph{Правая часть $\mathcal{F}_1$.}
\begin{equation}\label{eq:F1}
  \boxed{
    \mathcal{F}_1
    = \left[r\,\Lambda\,\frac{\pp p}{\pp r}\right]_{r_w}^{R_\ext}
    + Q_\text{quad}
    - I_{p_{wf}}\,\dot{p}_{wf},
  }
\end{equation}
где член с $\dot{p}_{wf}$ перенесён в правую часть,
т.\,к.\ $p_{wf}(t)$ --- заданная функция.

Граничные потоки:
\begin{equation}\label{eq:boundary_fluxes}
  \left[r\,\Lambda\,\frac{\pp p}{\pp r}\right]_{r_w}^{R_\ext}
  = \Lambda_\ext\,\frac{\Delta p_2}{L_e}
    - \underbrace{
      \Lambda_w\,\frac{\Delta p_1\,\gamma}
      {L_b}\,\lim_{\eta\to 0}\eta^{\gamma-1}
    }_{\to\,0\;\text{при}\;\gamma > 1;\;
      +\infty\;\text{при}\;\gamma < 1},
\end{equation}
где $\Lambda_\ext = \Lambda(p_0, T_0, 0, k)$ и
$\Lambda_w = \Lambda(p_{wf}, T_w, S_{g,w}, k)$ ---
точечные (не замороженные) значения на границах.
На контуре: $r\,\Lambda\,\pp p/\pp r\big|_{R_\ext}
= R_\ext\,\Lambda_\ext\,\Delta p_2/(R_\ext L_e)
= \Lambda_\ext\,\Delta p_2/L_e$.
На забое ($\eta = 0$): производная $\pp p/\pp r \sim \eta^{\gamma-1}$
расходится при $\gamma < 1$. Для baseline-параметров
$\gamma < 1$ всегда ($\Delta p_2/\Delta p_1 = 2/3$,
$L_b/L_e < 1$); фундаментальное ограничение параметризации
$\Phi(\eta) = \eta^\gamma$.

Квадратичный интеграл из~\eqref{eq:moment0}
редуцируется к зоне разгазирования, т.\,к.\ в однофазной зоне
$\pp R_s/\pp p = 0$ и $\pp R_s/\pp T = 0$ (фиксированный состав):
\begin{equation}\label{eq:Q_quad}
  Q_\text{quad}
  = \int_{r_w}^{r_b} r\,\frac{\lambda_o}{B_o}\!\left(
      \frac{\pp R_s}{\pp p}\,\frac{\pp p}{\pp r}
      + \frac{\pp R_s}{\pp T}\,\frac{\pp T}{\pp r}
    \right)\frac{\pp p}{\pp r}\,\dd r.
\end{equation}

Интеграл давления при нестационарном $p_{wf}$:
\begin{equation}\label{eq:I_pwf}
  I_{p_{wf}} = \int_{r_w}^{R_\ext} r\,C_p\,
    \frac{\pp p}{\pp p_{wf}}\,\dd r,
\end{equation}
при стационарном $p_{wf}$: $\dot{p}_{wf} = 0$ и вклад
обнуляется.

\emph{Знак:} $\mathcal{F}_1 < 0$ (отток через забой
доминирует над притоком через контур).
Это согласуется с $\mathcal{A}_{11} < 0$:
$\dot{r}_b = \mathcal{F}_1/\mathcal{A}_{11} > 0$,
т.\,е.\ зона разгазирования растёт.

\paragraph{Сингулярность при $\gamma < 1$.}
Подынтегральная $Q_\text{quad}$ содержит
$(\pp p/\pp r)^2 \sim \eta^{2(\gamma-1)}$. При $\gamma > 1/2$ ---
интегрируема; при $\gamma \leqslant 1/2$ --- расходится.
При начальном $r_b \approx r_w$ зона~1 вырождается
($L_b \to 0$), интеграл стремится к нулю.
При конечном $r_b$ с $\gamma < 1/2$: требуется
регуляризация (замена переменных или выделение
сингулярности при численном интегрировании).

\subsubsection{Коэффициенты второго уравнения (масса газа)}
\label{sssec:gas_mass_coefficients}

Подставляя~\eqref{eq:dp_dt}--\eqref{eq:dT_dt}
и~\eqref{eq:Sg_profile} в~\eqref{eq:Sg_moment0},
собираем коэффициенты при $\dot{r}_b$, $\dot{S}_{g,w}$,
$\dot{\Delta T}_w$.%
\footnote{FEAT-005b: коэффициенты уравнения массы газа.}

\emph{Область интегрирования.}
В~отличие от первого уравнения~(\S\ref{sssec:oil_mass_coefficients}),
где интеграл берётся по всей области $[r_w, R_\ext]$,
здесь все интегралы ограничены зоной разгазирования
$[r_w, r_b]$, т.\,к.\ $S_g \equiv 0$ при $r > r_b$.
Следовательно, достаточно одной замороженной зоны
($\bar{p}_1$, $\bar{T}_1$, $\bar{S}_{g,1}$,
см.~\S\ref{par:frozen_pvt}).

\emph{Вывод левой части.}
Применяя правило Лейбница к~$\mathcal{S}_0$~\eqref{eq:Sn}:
\begin{equation}\label{eq:dS0_dt}
  \frac{\dd\mathcal{S}_0}{\dd t}
  = \underbrace{r_b\,S_g(r_b)\,\dot{r}_b}_{=\,0}
    + \int_{r_w}^{r_b} r\,\frac{\pp S_g}{\pp t}\,\dd r,
\end{equation}
т.\,к.\ $S_g(r_b) = 0$. Поэтому $\dd\mathcal{S}_0/\dd t$
определяется только объёмным интегралом $\pp S_g/\pp t$.

\emph{Раскрытие $\pp S_g / \pp t$.}
Профиль~\eqref{eq:Sg_profile} зависит от $t$ через
$r_b(t)$ и $S_{g,w}(t)$:
\begin{equation}\label{eq:dSg_dt}
  \frac{\pp S_g}{\pp t}
  = \frac{\pp S_g}{\pp r_b}\,\dot{r}_b
    + \frac{\pp S_g}{\pp S_{g,w}}\,\dot{S}_{g,w},
\end{equation}
где
\begin{equation}\label{eq:dSg_drb}
  \frac{\pp S_g}{\pp r_b}
  = S_{g,w}\,\beta\,
    \frac{\eta^\beta}{r_b\,L_b},
  \qquad
  \frac{\pp S_g}{\pp S_{g,w}}
  = 1 - \eta^\beta,
\end{equation}
$\eta = \ln(r/r_w)/L_b$, $L_b = \ln(r_b/r_w)$.
Двойной минус: $\pp\eta/\pp r_b = -\eta/(r_b\,L_b)$
и $\pp S_g/\pp\eta = -S_{g,w}\,\beta\,\eta^{\beta-1}$
дают $\pp S_g / \pp r_b > 0$ (рост $r_b$ сдвигает фронт,
увеличивая $S_g$ внутри зоны).

\emph{Подстановка $\pp p / \pp t$ и $\pp T / \pp t$.}
В правой части~\eqref{eq:Sg_moment0} используем
разложения~\eqref{eq:dp_dt}--\eqref{eq:dT_dt},
учитывая $\pp T / \pp r_b = 0$~\eqref{eq:dT_drb}
и $\pp T/\pp \Delta T_w = -\Psi(\zeta)$~\eqref{eq:dT_dTw}.

\paragraph{Элемент $\mathcal{A}_{22}$.}
Коэффициент при $\dot{S}_{g,w}$. Из левой
части~\eqref{eq:Sg_moment0} (через~\eqref{eq:dSg_dt}):
\begin{equation}\label{eq:A22}
  \boxed{
    \mathcal{A}_{22}
    = \int_{r_w}^{r_b} r\,(1 - \eta^\beta)\,\dd r.
  }
\end{equation}
Геометрический интеграл, не зависящий от PVT.

\emph{Знак:} $\eta \in [0,1]$, поэтому $1-\eta^\beta > 0$
и $\mathcal{A}_{22} > 0$. При $r_b \to r_w$ зона
вырождается: $\mathcal{A}_{22} \to 0$.

\paragraph{Элемент $\mathcal{A}_{21}$.}
Коэффициент при $\dot{r}_b$. Левая часть
даёт интеграл~\eqref{eq:dSg_drb}, правая ---
вклад давления через $(\pp p/\pp r_b)\,\dot{r}_b$:
\begin{equation}\label{eq:A21}
  \boxed{
    \mathcal{A}_{21}
    = \frac{S_{g,w}\,\beta}{r_b\,L_b}\,
      \int_{r_w}^{r_b} r\,\eta^\beta\,\dd r
    - \int_{r_w}^{r_b} r\,S_o\,(c_r + c_o^\sat)\,
      \frac{\pp p}{\pp r_b}\,\dd r.
  }
\end{equation}
Первое слагаемое --- вклад геометрии
(движение фронта $r_b$), знак $> 0$.
Второе --- вклад давления:
$-\int r\,S_o\,(c_r+c_o^\sat)\,(\pp p/\pp r_b)\,\dd r$.
Знак этого слагаемого зависит от $c_r + c_o^\sat$.
При типичных условиях $|c_o^\sat| \gg c_r$
(поскольку $c_o^\sat < 0$~\eqref{eq:co_sat}),
и $c_r + c_o^\sat < 0$, что даёт второму слагаемому
знак $< 0$. Однако геометрический вклад (первое
слагаемое) доминирует, поэтому $\mathcal{A}_{21} > 0$.

Температурный вклад в $\mathcal{A}_{21}$ обращается
в нуль, т.\,к.\ $\pp T / \pp r_b = 0$~\eqref{eq:dT_drb}.

PVT-свойства замораживаются по зоне~1:
$\bar{S}_{o,1} = 1 - \bar{S}_{g,1}$,
$\bar{c}_{o,1}^\sat = c_o^\sat(\bar{p}_1, \bar{T}_1)$.

\paragraph{Элемент $\mathcal{A}_{23}$.}
Коэффициент при $\dot{\Delta T}_w$. Из правой
части~\eqref{eq:Sg_moment0} --- вклад температуры
через $\pp T / \pp \Delta T_w = -\Psi(\zeta)$~\eqref{eq:dT_dTw}:
\begin{equation}\label{eq:A23}
  \boxed{
    \mathcal{A}_{23}
    = -\int_{r_w}^{r_b} r\,S_o\,\beta_o^\sat\,\Psi(\zeta)\,\dd r.
  }
\end{equation}
\emph{Знак.}
$\Psi(\zeta) > 0$, $S_o > 0$, поэтому
$\mathrm{sign}(\mathcal{A}_{23}) = -\mathrm{sign}(\beta_o^\sat)$.
Величина $\beta_o^\sat$~\eqref{eq:beta_o_sat}:
$\beta_o^\sat = \beta_o + (1/B_o)\,
(\pp B_o/\pp R_s)_p\,(\pp R_s/\pp T)_p$.
Тепловое расширение $\beta_o > 0$, но коррекция
$(1/B_o)(\pp B_o / \pp R_s)(\pp R_s / \pp T)$ отрицательна
(т.\,к.\ $\pp R_s / \pp T < 0$ при типичных условиях),
поэтому знак $\beta_o^\sat$ зависит от конкуренции
двух вкладов. При baseline-условиях
($p \approx 16{,}5$~МПа, $T = 348$~К)
коррекция превышает по модулю $\beta_o$, и
$\beta_o^\sat < 0$,
поэтому $\mathcal{A}_{23} > 0$.

\emph{Вывод.}
В~\eqref{eq:Sg_moment0} температурный член
$-\int r\,S_o\,\beta_o^\sat\,(\pp T/\pp t)\,\dd r$
стоит в правой части. Подставляя
$\pp T / \pp t = (\pp T / \pp\Delta T_w)\,\dot{\Delta T}_w
= -\Psi\,\dot{\Delta T}_w$ и перенося $\dot{\Delta T}_w$
в левую часть:
\[
  \mathcal{A}_{23} = -(-1)\!\int r\,S_o\,\beta_o^\sat\,\Psi\,\dd r
  \;\Longrightarrow\;
  \mathcal{A}_{23} = +\!\int r\,S_o\,\beta_o^\sat\,\Psi\,\dd r?
\]
Нет: минус перед интегралом $-\int r\,S_o\,\beta_o^\sat$
в~\eqref{eq:Sg_moment0} и минус от $\pp T/\pp\Delta T_w = -\Psi$
дают при переносе в левую часть:
\[
  \text{ЛЧ} \leftarrow -\bigl[-S_o\,\beta_o^\sat\,(-\Psi)\bigr]\,\dot{\Delta T}_w
  = -S_o\,\beta_o^\sat\,\Psi\,\dot{\Delta T}_w,
\]
что и даёт~\eqref{eq:A23}.

\paragraph{Правая часть $\mathcal{F}_2$.}
Из граничных потоков~\eqref{eq:Sg_flux}
(с замороженным $B_o$ из зоны~1):
\begin{equation}\label{eq:F2}
  \boxed{
    \mathcal{F}_2
    = -\frac{B_o}{\phi}\!\left[
        \frac{r\,\lambda_o}{B_o}\,\frac{\pp p}{\pp r}
      \right]_{r_w}^{r_b}
    + I_{p_{wf}}^{(2)}\,\dot{p}_{wf},
  }
\end{equation}
где интеграл давления при нестационарном $p_{wf}$:
\begin{equation}\label{eq:I_pwf_2}
  I_{p_{wf}}^{(2)} = \int_{r_w}^{r_b} r\,S_o\,(c_r+c_o^\sat)\,
    \frac{\pp p}{\pp p_{wf}}\,\dd r.
\end{equation}

\emph{Структурное отличие знака от $\mathcal{F}_1$.}
В первом уравнении~\eqref{eq:moment0} интеграл давления
$\int r\,C_p\,(\pp p/\pp t)\,\dd r$ стоит в левой части.
При выделении $\dot{p}_{wf}$ этот интеграл переносится
в правую часть с~переменой знака: $-I_{p_{wf}}\,\dot{p}_{wf}$
в~\eqref{eq:F1}. Во~втором уравнении~\eqref{eq:Sg_moment0}
интеграл давления $\int r\,S_o\,(c_r+c_o^\sat)\,(\pp p/\pp t)
\,\dd r$ стоит в правой части, поэтому
$I_{p_{wf}}^{(2)}\,\dot{p}_{wf}$ остаётся с~\emph{тем~же}
знаком: $+I_{p_{wf}}^{(2)}\,\dot{p}_{wf}$ в~\eqref{eq:F2}.

\emph{Знак:} при стационарном $p_{wf}$ ($\dot{p}_{wf} = 0$)
правая часть $\mathcal{F}_2$ определяется только граничным
потоком. При $\gamma < 1$ поток через забой ($r = r_w$)
доминирует: $\mathcal{F}_2 > 0$.

\paragraph{Сингулярность потока при $\gamma < 1$.}
Аналогично $\mathcal{F}_1$~\eqref{eq:boundary_fluxes},
поток $r\,(\lambda_o/B_o)\,(\pp p/\pp r)$ при $r \to r_w$
содержит $\eta^{\gamma-1}$, расходящийся при $\gamma < 1$.
В коде применяется $\varepsilon$-отступ:
$r_w + \varepsilon$, $\varepsilon = 10^{-3}\,r_w$.


\subsubsection{Коэффициенты третьего уравнения (энергия)}

Нулевой момент уравнения энергии~\eqref{eq:T_moment0}
содержит ёмкостный, конвективный и источниковые члены.
Подставляя~\eqref{eq:dT_dt} и раскрывая
$\partial p/\partial t$~\eqref{eq:dp_dt},
получаем строку $(\mathcal{A}_{31},\, 0,\, \mathcal{A}_{33})$
и правую часть $\mathcal{F}_3$.

Элемент $\mathcal{A}_{32} = 0$ по расщеплению~\eqref{eq:matrix_structure}:
температурный профиль $\Psi(\zeta)$ не зависит от~$S_{g,w}$.

\paragraph{Элемент $\mathcal{A}_{31}$.}

Вклад даёт только $Q_{\text{ad}}^*$~\eqref{eq:Qad_star},
содержащий $\partial p/\partial t$.
Ёмкостный член~(I) не вносит вклада:
$\partial T/\partial r_b = 0$~\eqref{eq:dT_drb}.
Конвективный~(II) содержит $\partial T/\partial r$, а не
$\partial T/\partial t$, и не зависит от~$\dot{r}_b$.

Коэффициент при $\dot{r}_b$ стоит в правой части (eq:T\_moment0),
поэтому при переносе в матрицу ($A\dot{q} = F$) знак меняется:
\begin{equation}\label{eq:A31}
  \boxed{
    \mathcal{A}_{31}
    = -\int_{r_w}^{R_{\text{ext}}} r\,\bar{Q}_{\text{ad}}^{*}\,
      \frac{\partial p}{\partial r_b}\,dr
  }
\end{equation}
где $\bar{Q}_{\text{ad}}^{*}$ --- коэффициент при $\partial p/\partial t$
в формуле $Q_{\text{ad}}^*$:
\begin{equation}\label{eq:Qad_star_coeff}
  \bar{Q}_{\text{ad}}^{*}(p,T,S_g)
  = \phi\!\left[T\sum_\alpha S_\alpha\,\alpha_\alpha
    + c_r\bigl(p + (\rho c_p)_s\,T\bigr)\right]
  + L\,\rho_g^{\text{ст}}\,\phi\,\xi_o\,\frac{\partial R_s}{\partial p}.
\end{equation}

Двухзонное разбиение:
\begin{equation}\label{eq:A31_zones}
  \mathcal{A}_{31}
  = -\bar{Q}_{\text{ad},1}^{*}
    \int_{r_w}^{r_b} r\,\frac{\partial p}{\partial r_b}\bigg|_{\text{inner}} dr
    - \bar{Q}_{\text{ad},2}^{*}
    \int_{r_b}^{R_{\text{ext}}} r\,\frac{\partial p}{\partial r_b}\bigg|_{\text{outer}} dr.
\end{equation}

\emph{Знак:} $\bar{Q}_{\text{ad}}^{*} > 0$,
$\partial p/\partial r_b < 0$, интеграл $< 0$,
минус от переноса RHS$\to$LHS
$\Rightarrow$ $\mathcal{A}_{31} > 0$.

\emph{Структура:} аналогична $\mathcal{A}_{11}$~\eqref{eq:A11_zones}
с заменой $C_p \to \bar{Q}_{\text{ad}}^{*}$ и противоположным знаком.

\paragraph{Элемент $\mathcal{A}_{33}$.}

Вклад даёт ёмкостный член~(I) через
$\partial T/\partial \Delta T_w = -\Psi(\zeta)$~\eqref{eq:dT_dTw}.
Конвективный и источниковые не зависят от~$\dot{\Delta T}_w$.

\begin{equation}\label{eq:A33}
  \boxed{
    \mathcal{A}_{33}
    = -\int_{r_w}^{R_{\text{ext}}} r\,(\rho c_p)_{\text{eff}}^{*}\,
      \Psi(\zeta)\,dr
  }
\end{equation}

Двухзонное разбиение:
\begin{equation}\label{eq:A33_zones}
  \mathcal{A}_{33}
  = -(\rho c_p)_{\text{eff},1}^{*}
    \int_{r_w}^{r_b} r\,\Psi(\zeta)\,dr
    - (\rho c_p)_{\text{eff},2}^{*}
    \int_{r_b}^{R_{\text{ext}}} r\,\Psi(\zeta)\,dr.
\end{equation}

\emph{Знак:}
$(\rho c_p)_{\text{eff}}^{*} > 0$, $\Psi > 0$
$\Rightarrow$ $\mathcal{A}_{33} < 0$.
Минус от $\partial T/\partial\Delta T_w = -\Psi$.

\emph{Структура:} аналогична $\mathcal{A}_{13}$~\eqref{eq:A13}
с заменой $C_T \to (\rho c_p)_{\text{eff}}^{*}$.

\paragraph{Правая часть $\mathcal{F}_3$.}

Собирает все члены, не зависящие от
$\dot{r}_b$, $\dot{S}_{g,w}$, $\dot{\Delta T}_w$:
\begin{equation}\label{eq:F3}
  \boxed{
    \mathcal{F}_3
    = \underbrace{R_{\text{ext}}\,\lambda_{\text{eff}}\,T'(R_{\text{ext}})}_{
        \text{теплопроводность}}
    + \underbrace{\int_{r_w}^{R_{\text{ext}}} r\,\bar{C}_{\text{conv}}^{*}\,
      \frac{\partial p}{\partial r}\,\frac{\partial T}{\partial r}\,dr}_{
        \text{конвекция}}
    + \underbrace{\int_{r_w}^{R_{\text{ext}}} r\,\bar{Q}_{\text{JT}}^{*}\,
      \left(\frac{\partial p}{\partial r}\right)^{\!2} dr}_{
        \text{JT}}
    + \underbrace{I_{p_{wf}}^{(3)}\,\dot{p}_{wf}}_{
        \text{нестац.\ }p_{wf}}
  }
\end{equation}

Замороженные коэффициенты:
\begin{align}
  \bar{C}_{\text{conv}}^{*}
    &= \rho_o\,c_{p,o}\,\lambda_o + \rho_g\,c_{p,g}\,\lambda_g
       - L\,\rho_g^{\text{ст}}\,\frac{\lambda_o}{B_o}\,
         \frac{\partial R_s}{\partial T},
  \label{eq:Cconv_star}\\[4pt]
  \bar{Q}_{\text{JT}}^{*}
    &= \sum_\alpha \lambda_\alpha\,(1 - \alpha_\alpha\,T)
       - L\,\rho_g^{\text{ст}}\,\frac{\lambda_o}{B_o}\,
         \frac{\partial R_s}{\partial p},
  \label{eq:QJT_star_coeff}\\[4pt]
  I_{p_{wf}}^{(3)}
    &= \int_{r_w}^{R_{\text{ext}}} r\,\bar{Q}_{\text{ad}}^{*}\,
       \frac{\partial p}{\partial p_{wf}}\,dr.
  \label{eq:I_pwf_energy}
\end{align}

\emph{Знаки конвективного и JT интегралов.}
Конвективный: два минуса (перенос LHS$\to$RHS
и Дарси $u = -\lambda\,\partial p/\partial r$) компенсируются
$\Rightarrow$ знак~<<$+$>>.
JT: $Q_{\text{JT}}^* \propto (\partial p/\partial r)^2 \geq 0$,
стоит в RHS --- знак~<<$+$>>.

\emph{Замечание:} знак $\bar{Q}_{\text{JT}}^{*}$ не фиксирован.
При сильной дегазации ($L\,\rho_g^{\text{ст}}\,(\lambda_o/B_o)\,
\partial R_s/\partial p$ велико) он может быть отрицательным ---
дегазационное охлаждение доминирует над JT-нагревом.

\paragraph{Сингулярности.}
JT-интеграл содержит $(\partial p/\partial r)^2 \sim \eta^{2(\gamma-1)}$,
интегрируемый при $\gamma > 1/2$ (как в $Q_{\text{quad}}$
из~$\mathcal{F}_1$).
Конвективный содержит произведение
$(\partial p/\partial r)(\partial T/\partial r)$;
при $r \to r_w$: $\partial T/\partial r = 0$ (ГУ),
поэтому произведение конечно при $\gamma + \sigma > 1$
(baseline: $\gamma + \sigma \approx 2.6$).

\paragraph{Проверка размерностей.}
Все члены $A\dot{q} = F$ имеют размерность [Вт/м]:
\begin{itemize}
  \item $[\mathcal{A}_{31}] = [\text{Дж/м}^2]$;
        $[\mathcal{A}_{31}\,\dot{r}_b] = [\text{Вт/м}]$ \checkmark
  \item $[\mathcal{A}_{33}] = [\text{Дж/(м\textperiodcentered К)}]$;
        $[\mathcal{A}_{33}\,\dot{\Delta T}_w] = [\text{Вт/м}]$ \checkmark
  \item $[\mathcal{F}_3] = [\text{Вт/м}]$ \checkmark
\end{itemize}


%=============================================================================
\subsection{Расширение ODE: показатель формы $\beta(t)$}
\label{ssec:beta_ode}
%=============================================================================

\paragraph{Мотивация.}
При фиксированном $\beta = \text{const}$ система~\eqref{eq:ode_system}
из трёх~ODE ($r_b$, $S_{g,w}$, $\Delta T_w$) порождает
затухающие осцилляции всех переменных при нестационарных
граничных условиях. Газонасыщенность $S_{g,w}$ может уходить
в отрицательные значения, что нефизично.

Причина: три степени свободы при фиксированных формах профилей
не могут описать перестройку пространственного профиля $S_g(r)$.
Добавление $\beta(t)$ как четвёртой ODE-переменной расширяет
пространство допустимых профилей.

\paragraph{Расширение разложения $\pp S_g/\pp t$.}
При $\beta = \beta(t)$ разложение~\eqref{eq:dSg_dt}
дополняется третьим членом:
\begin{equation}\label{eq:dSg_dt_ext}
  \frac{\pp S_g}{\pp t}
  = \frac{\pp S_g}{\pp r_b}\,\dot{r}_b
    + \frac{\pp S_g}{\pp S_{g,w}}\,\dot{S}_{g,w}
    + \frac{\pp S_g}{\pp \beta}\,\dot{\beta},
\end{equation}
где
\begin{equation}\label{eq:dSg_dbeta}
  \frac{\pp S_g}{\pp \beta}
  = -S_{g,w}\,\eta^\beta\,\ln\eta.
\end{equation}

\emph{Знак:} $\ln\eta < 0$ при $\eta \in (0,1)$,
поэтому $\pp S_g/\pp\beta > 0$:
увеличение $\beta$ увеличивает $S_g$ (профиль
становится более <<плоским>>, приближаясь к $S_{g,w}$).

\paragraph{Элемент $\mathcal{A}_{2\beta}$ (при $\dot\beta$
во~2-м уравнении).}
Подставляя~\eqref{eq:dSg_dt_ext}
в нулевой момент~\eqref{eq:Sg_moment0},
получаем новый коэффициент при $\dot\beta$:
\begin{equation}\label{eq:A2beta}
  \boxed{
    \mathcal{A}_{2\beta}
    = \int_{r_w}^{r_b} r\,
      \frac{\pp S_g}{\pp\beta}\,\dd r
    = -S_{g,w}\!\int_{r_w}^{r_b}
      r\,\eta^\beta\,\ln\eta\,\dd r.
  }
\end{equation}
\emph{Знак:} $\ln\eta < 0$, минус перед интегралом,
поэтому $\mathcal{A}_{2\beta} > 0$ при $S_{g,w} > 0$.

Уравнение~\eqref{eq:ode2} расширяется:
\begin{equation}\label{eq:ode2_ext}
  \mathcal{A}_{21}\,\dot{r}_b
  + \mathcal{A}_{22}\,\dot{S}_{g,w}
  + \mathcal{A}_{2\beta}\,\dot{\beta}
  + \mathcal{A}_{24}\,\dot{\Delta T}_w
  = \mathcal{F}_2,
\end{equation}
где $\mathcal{A}_{24} \equiv \mathcal{A}_{23}$
из~\eqref{eq:A23} (перенумерация).

\paragraph{Выбор весовой функции 4-го момента.}
Для получения уравнения на $\dot\beta$ нужен
четвёртый интегральный момент уравнения насыщенности.
Рассмотрим вес $w_4(r) = r\,\eta^\beta$. Он:
\begin{itemize}[nosep]
  \item даёт замкнутые интегралы (степенные по $\eta$);
  \item подавляет сингулярность $\pp p/\pp r$ при $r \to r_w$
    ($\eta^\beta \to 0$ при $\eta \to 0$, $\beta > 0$);
  \item линейно независим от веса нулевого момента
    (вес~$r$): $\int_0^1 (1-\eta^\beta)\,\eta^\beta\,\dd\eta
    \ne 0$.
\end{itemize}

\paragraph{Четвёртый момент.}
Умножаем~\eqref{eq:Sg_eq2} на $r\,\eta^\beta$
и интегрируем от $r_w$ до $r_b$:
\begin{equation}\label{eq:Sg_moment_beta}
  \int_{r_w}^{r_b} r\,\eta^\beta\,
    \frac{\pp S_g}{\pp t}\,\dd r
  = \int_{r_w}^{r_b} r\,\eta^\beta
    \left[
      S_o\,(c_r + c_o^\sat)\,\frac{\pp p}{\pp t}
      - S_o\,\beta_o^\sat\,\frac{\pp T}{\pp t}
    \right]\dd r
  + F_4^\div,
\end{equation}
где $F_4^\div$ --- вклад дивергентного члена
(см.\ ниже, \eqref{eq:F4_boundary}--\eqref{eq:F4_volume}).

\paragraph{Левая часть.}
Подставляя параметризацию
$S_g = S_{g,w}(1 - \eta^\beta)$,
$\pp S_g/\pp t$ содержит
$\dot{r}_b$, $\dot{S}_{g,w}$, $\dot\beta$
(из~\eqref{eq:dSg_dt_ext}).
Интеграл $\int r\,\eta^\beta\,(\pp S_g/\pp t)\,\dd r$
тождественно равен
\begin{equation}\label{eq:dM4_dt}
  \int_{r_w}^{r_b} r\,\eta^\beta\,
    \frac{\pp S_g}{\pp t}\,\dd r
  = \int_{r_w}^{r_b} r\,\eta^\beta\,
    \biggl[
      \frac{\pp S_g}{\pp r_b}\,\dot{r}_b
      + \frac{\pp S_g}{\pp S_{g,w}}\,\dot{S}_{g,w}
      + \frac{\pp S_g}{\pp\beta}\,\dot\beta
    \biggr]\,\dd r.
\end{equation}
Подставляя производные~\eqref{eq:dSg_drb}
и~\eqref{eq:dSg_dbeta}:
\begin{align}
  \eta^\beta\,\frac{\pp S_g}{\pp r_b}
  &= S_{g,w}\,\beta\,
     \frac{\eta^{2\beta}}{r_b\,L_b},
  \label{eq:eta_dSg_drb} \\
  \eta^\beta\,\frac{\pp S_g}{\pp S_{g,w}}
  &= \eta^\beta\,(1-\eta^\beta),
  \label{eq:eta_dSg_dSgw} \\
  \eta^\beta\,\frac{\pp S_g}{\pp\beta}
  &= -S_{g,w}\,\eta^{2\beta}\,\ln\eta.
  \label{eq:eta_dSg_dbeta}
\end{align}

\paragraph{Элемент $\mathcal{A}_{41}$ (при $\dot{r}_b$).}
Из левой части~\eqref{eq:dM4_dt}
(через~\eqref{eq:eta_dSg_drb})
и из правой части~\eqref{eq:Sg_moment_beta}
(давление через $(\pp p/\pp r_b)\dot{r}_b$,
аналогично $\mathcal{A}_{21}$~\eqref{eq:A21}
но с весом $\eta^\beta$):
\begin{equation}\label{eq:A41}
  \boxed{
    \mathcal{A}_{41}
    = \frac{S_{g,w}\,\beta}{r_b\,L_b}
      \int_{r_w}^{r_b} r\,\eta^{2\beta}\,\dd r
    - \int_{r_w}^{r_b} r\,\eta^\beta\,
      S_o\,(c_r + c_o^\sat)\,
      \frac{\pp p}{\pp r_b}\,\dd r.
  }
\end{equation}
Структура идентична $\mathcal{A}_{21}$~\eqref{eq:A21}
с заменой $\eta^\beta \to \eta^{2\beta}$ в геометрическом
интеграле и добавлением веса $\eta^\beta$
во~втором интеграле.

\paragraph{Элемент $\mathcal{A}_{42}$ (при $\dot{S}_{g,w}$).}
Из~\eqref{eq:eta_dSg_dSgw}:
\begin{equation}\label{eq:A42}
  \boxed{
    \mathcal{A}_{42}
    = \int_{r_w}^{r_b} r\,\eta^\beta\,(1 - \eta^\beta)\,\dd r.
  }
\end{equation}
Геометрический, $\mathcal{A}_{42} > 0$.

\paragraph{Элемент $\mathcal{A}_{43}$ (при $\dot\beta$).}
Из~\eqref{eq:eta_dSg_dbeta}:
\begin{equation}\label{eq:A43}
  \boxed{
    \mathcal{A}_{43}
    = -S_{g,w}\!\int_{r_w}^{r_b} r\,
      \eta^{2\beta}\,\ln\eta\,\dd r.
  }
\end{equation}

\emph{Знак.} $\ln\eta < 0$ при $\eta \in (0,1)$,
минус перед интегралом, поэтому
$\mathcal{A}_{43} > 0$ при $S_{g,w} > 0$.

\paragraph{Элемент $\mathcal{A}_{44}$
  (при $\dot{\Delta T}_w$).}
Из правой части (температура через
$\pp T/\pp\Delta T_w = -\Psi$,
аналогично $\mathcal{A}_{23}$~\eqref{eq:A23}
но с весом $\eta^\beta$):
\begin{equation}\label{eq:A44}
  \boxed{
    \mathcal{A}_{44}
    = -\int_{r_w}^{r_b} r\,\eta^\beta\,
      S_o\,\beta_o^\sat\,\Psi(\zeta)\,\dd r.
  }
\end{equation}

\paragraph{Правая часть $\mathcal{F}_4$.}
Дивергентный член интегрируется с весом $r\,\eta^\beta$.
Интегрирование по частям:
\begin{equation}\label{eq:F4_boundary}
  F_4^\div
  = -\frac{B_o}{\phi}
    \biggl[
      \underbrace{\eta^\beta\,
        \frac{r\,\lambda_o}{B_o}\,p'}_{r_w}^{r_b}
    \biggr]
  + \frac{B_o}{\phi}
    \int_{r_w}^{r_b}
    \frac{\dd\eta^\beta}{\dd r}\,
    \frac{r\,\lambda_o}{B_o}\,p'\,\dd r.
\end{equation}

Граничные значения: при $r = r_b$: $\eta^\beta = 1$
(поток на фронте, как в $\mathcal{F}_2$).
При $r = r_w$: $\eta^\beta = 0$
--- \textbf{сингулярность $p'$ подавлена}.

Объёмный интеграл: $\dd\eta^\beta/\dd r
= \beta\,\eta^{\beta-1}/(r\,L_b)$, поэтому
\begin{equation}\label{eq:F4_volume}
  F_4^{(V)}
  = \frac{B_o}{\phi}\,\frac{\beta}{L_b}
    \int_{r_w}^{r_b}
    \eta^{\beta-1}\,
    \frac{\lambda_o}{B_o}\,p'\,\dd r.
\end{equation}

С замороженным $\lambda_o/B_o$ и профилем
давления $p' \propto \eta^{\gamma-1}/(r\,L_b)$
подынтегральное $\propto \eta^{\beta+\gamma-2}$;
интеграл сходится при $\beta + \gamma > 1$
(всегда: $\beta \geqslant 1$, $\gamma > 0$).

Итого:
\begin{equation}\label{eq:F4}
  \boxed{
    \mathcal{F}_4
    = -\frac{B_o}{\phi}
      \left.\frac{r\,\lambda_o}{B_o}\,p'
      \right|_{r=r_b}
    + F_4^{(V)}
    + I_{p_{wf}}^{(4)}\,\dot{p}_{wf},
  }
\end{equation}
где $I_{p_{wf}}^{(4)}$ --- вклад нестационарного $p_{wf}$
(аналогичен $I_{p_{wf}}^{(2)}$~\eqref{eq:I_pwf_2}
с весом $\eta^\beta$).

\emph{Отличие от $\mathcal{F}_2$:}
в $\mathcal{F}_2$~\eqref{eq:F2} поток на~$r_w$
содержит сингулярность при $\gamma < 1$.
В $\mathcal{F}_4$ поток на~$r_w$ обнуляется
весом $\eta^\beta(r_w) = 0$.

\paragraph{Новые замкнутые интегралы.}
Базовый интеграл с логарифмом:
\begin{equation}\label{eq:K_base}
  \int_0^1 \eta^m\,\ln\eta\,\dd\eta = -\frac{1}{(m{+}1)^2},
  \qquad m > -1
\end{equation}
(интегрирование по частям: $u = \ln\eta$,
$\dd v = \eta^m\,\dd\eta$).

Для $\mathcal{A}_{42}$:
\begin{equation}\label{eq:I_S_tilde}
  \widetilde{I}_S^{(n)}(\beta)
  = \int_0^1 \eta^{n+\beta}\,(1-\eta^\beta)\,\dd\eta
  = \frac{1}{n+\beta+1} - \frac{1}{n+2\beta+1}.
\end{equation}

Для $\mathcal{A}_{43}$ (при интегрировании по $\eta$):
\begin{equation}\label{eq:K_S}
  \int_0^1 \eta^{2\beta}\,\ln\eta\,\dd\eta
  = -\frac{1}{(2\beta{+}1)^2}.
\end{equation}

Для $\mathcal{A}_{2\beta}$~\eqref{eq:A2beta}:
\begin{equation}\label{eq:K_A2beta}
  \int_0^1 \eta^{\beta}\,\ln\eta\,\dd\eta
  = -\frac{1}{(\beta{+}1)^2}.
\end{equation}

\emph{Замечание.} Интегралы записаны в~переменных
$\eta \in [0,1]$.
При вычислении $\mathcal{A}_{4j}$ через \texttt{quad}
по~$r \in [r_w, r_b]$ множитель $r$ и якобиан
$\dd r = r\,L_b\,\dd\eta$ учитываются автоматически;
замкнутые формулы выше служат для оценок
и предельных случаев.

\paragraph{Расширенная матрица 4$\times$4.}
С учётом переменного $\beta(t)$,
вектор неизвестных:
$\mathbf{q} = (r_b,\,S_{g,w},\,\beta,\,\Delta T_w)^\top$.
Система:
\begin{equation}\label{eq:ode_system_4x4}
  \boxed{
    \begin{pmatrix}
      \mathcal{A}_{11} & 0
        & 0 & \mathcal{A}_{14} \\
      \mathcal{A}_{21} & \mathcal{A}_{22}
        & \mathcal{A}_{2\beta} & \mathcal{A}_{24} \\
      \mathcal{A}_{31} & 0
        & 0 & \mathcal{A}_{34} \\
      \mathcal{A}_{41} & \mathcal{A}_{42}
        & \mathcal{A}_{43} & \mathcal{A}_{44}
    \end{pmatrix}
    \begin{pmatrix}
      \dot{r}_b \\ \dot{S}_{g,w} \\ \dot\beta
      \\ \dot{\Delta T}_w
    \end{pmatrix}
    =
    \begin{pmatrix}
      \mathcal{F}_1 \\ \mathcal{F}_2
      \\ \mathcal{F}_3 \\ \mathcal{F}_4
    \end{pmatrix}.
  }
\end{equation}

Строки~1 и~3 (давление и энергия) не содержат $\dot\beta$:
профили $p(r)$ и $T(r)$ не зависят от $\beta$
(при замороженных PVT: $C_p$, $C_T$ вычисляются
при среднем $\bar{S}_g = \tfrac{1}{2}S_{g,w}$,
не зависящем от $\beta$).

$\mathcal{A}_{14} \equiv \mathcal{A}_{13}$~\eqref{eq:A13},
$\mathcal{A}_{24} \equiv \mathcal{A}_{23}$~\eqref{eq:A23},
$\mathcal{A}_{34} \equiv \mathcal{A}_{33}$~\eqref{eq:A33}
(перенумерация: $\Delta T_w$ переходит в~4-й столбец).

\paragraph{Допущение: замороженные PVT и $\beta$.}
PVT-свойства ($\bar{S}_o$, $\bar{c}_o^\sat$,
$\bar\beta_o^\sat$) вычисляются при среднем
$\bar{S}_g = \tfrac{1}{2}S_{g,w}$, не зависящем от~$\beta$.
Точное среднее $\langle S_g\rangle = S_{g,w}(1 - 1/(\beta{+}1))$
зависит от~$\beta$, но различие при $\beta \in [1,\,4]$
составляет $< 20\%$ и поглощается приближением
замороженных коэффициентов.

\paragraph{Предельные случаи.}

\emph{$\beta = \mathrm{const}$ ($\dot\beta = 0$).}
Четвёртое уравнение становится алгебраическим ограничением;
первые три совпадают с~\eqref{eq:ode_system}
(столбец~$\dot\beta$ зануляется).
При $\beta = \mathrm{const}$ 4-ю строку отбрасываем,
восстанавливая исходную систему~3$\times$3.

\emph{$S_{g,w} \to 0$.}
$\mathcal{A}_{43} \propto S_{g,w} \to 0$;
$\mathcal{A}_{2\beta} \propto S_{g,w} \to 0$.
Уравнение для $\dot\beta$ вырождается:
форма пустого профиля неопределена.
Рекомендуется \verb|clamp|: $\dot\beta = 0$
при $S_{g,w} < \varepsilon$.

\emph{$\beta \to 1$.}
Профиль $S_g = S_{g,w}(1-\eta)$ (линейный).
Все $\mathcal{A}_{4j}$, $\mathcal{F}_4$ конечны
($\widetilde{I}_S^{(0)} = 1/2 - 1/3 = 1/6$,
$\int_0^1\eta^2\ln\eta\,\dd\eta = -1/9$).

\emph{$\beta \to \infty$.}
$S_g \to S_{g,w}$ (ступенька),
$\widetilde{I}_S^{(n)} \to 0$,
$\int_0^1\eta^{2\beta}\ln\eta\,\dd\eta \to 0$,
$\mathcal{A}_{42} \to 0$,
$\mathcal{A}_{43} \to 0$.
Уравнение для $\dot\beta$ стабилизирует $\beta$.

\paragraph{Таблица: формула $\to$ код (для FEAT-010).}

\begin{center}
\small
\begin{tabular}{lll}
  \textbf{Формула} & \textbf{Python-функция} & \textbf{Модуль} \\
  \hline
  \eqref{eq:A2beta} & \texttt{gas\_A2beta()} & \texttt{coefficients.py} \\
  \eqref{eq:A41}    & \texttt{gas\_A41()}    & \texttt{coefficients.py} \\
  \eqref{eq:A42}    & \texttt{gas\_A42()}    & \texttt{coefficients.py} \\
  \eqref{eq:A43}    & \texttt{gas\_A43()}    & \texttt{coefficients.py} \\
  \eqref{eq:A44}    & \texttt{gas\_A44()}    & \texttt{coefficients.py} \\
  \eqref{eq:F4}     & \texttt{gas\_F4()}     & \texttt{coefficients.py} \\
  \eqref{eq:I_S_tilde} & \texttt{I\_S\_tilde()} & \texttt{integrals.py} \\
  \eqref{eq:K_S}       & \texttt{K\_S()}        & \texttt{integrals.py} \\
\end{tabular}
\end{center}


%=============================================================================
\section{Начальная стадия: автомодельная асимптотика при $t \to 0$}
\label{sec:initial_stage}
%=============================================================================

При $t = 0$ давление скачком снижается до $p_{wf} < p_b$.
Формально $r_b(0) = r_w$, $S_{g,w}(0) = 0$, $\Delta T_w(0) = 0$.
При $t = 0^+$ $S_{g,w}$ скачком принимает конечное
значение, но $L_b = 0$, $\gamma = 0$, и матрица
системы~\eqref{eq:ode_system} вырождается.

\subsection{Физическая картина}

При $t < 0$ давление $p = p_0 > p_b$, температура $T = T_0$:
газ растворён, фильтрация однофазная.
При $t = 0$ мгновенно $p(r_w) = p_{wf} < p_b$.

На малых временах $t \ll r_w^2/\chi_0$ задача описывается
однофазной диффузией давления. Температура практически
не меняется: характерное время теплового отклика
$\tau_T \sim r_w^2/a$ ($a = \lambda_\eff/(\rho c_p)_\eff$
--- температуропроводность), и при типичных параметрах
$a \sim 10^{-7}$~м$^2$/с, $r_w = 0{,}1$~м:
$\tau_T \sim 10^5$~с, т.\,е.\
$\tau_T/\tau_\mathrm{rel} \sim 10^7$.

Таким образом, на стадии зарождения зоны разгазирования
температурное поле <<заморожено>>: $T \approx T_0$,
$\Delta T_w \approx 0$. Начальные условия для $r_b$
и $S_{g,w}$ определяются из изотермической задачи.

\subsection{Локальная задача диффузии}

При $\delta = r_b - r_w \ll r_w$:
\begin{equation}\label{eq:diffusion}
  \frac{\pp p}{\pp t} = \chi_0\,\frac{\pp^2 p}{\pp \xi^2},
  \qquad
  \chi_0 = \frac{\Lambda_0}{C_p^{(0)}},
  \qquad
  \xi = r - r_w,
\end{equation}
где $\Lambda_0 = k/(B_g\,\mu_o)$ при $S_g = 0$,
$C_p^{(0)}$ --- коэффициент~\eqref{eq:Cp}
при $S_g = 0$, $T = T_0$.
Член $C_T\,\dot{T}$ в~\eqref{eq:p_eq} обнуляется не потому,
что $C_T = 0$, а потому, что на этой стадии $\dot{T} \approx 0$
($\tau_T \gg \tau_\mathrm{rel}$).

Решение:
\begin{equation}\label{eq:p_erfc}
  p(\xi, t) = p_{wf} + (p_0 - p_{wf})\,
    \operatorname{erf}\!\left(\frac{\xi}{2\sqrt{\chi_0\,t}}\right).
\end{equation}

\subsection{Скорость роста границы разгазирования}

Условие $p(r_b) = p_b$ даёт:
\begin{equation}\label{eq:delta_sqrt}
  \boxed{
    \delta(t) = r_b(t) - r_w = 2\,\alpha_*\,\sqrt{\chi_0\,t},
    \qquad
    \alpha_* = \operatorname{erf}^{-1}\!\left(
      \frac{p_b - p_{wf}}{p_0 - p_{wf}}
    \right).
  }
\end{equation}

В отличие от задачи Стефана, $r_b$ --- кинематически
определённая изобара на гладком erfc-профиле.

\subsection{Асимптотика параметров профилей}

При $\delta \ll r_w$:
\begin{equation}\label{eq:Lb_asymp}
  L_b \approx \frac{\delta}{r_w},
  \qquad
  L_e \to L_0 = \ln\!\frac{R_\ext}{r_w}.
\end{equation}
\begin{equation}\label{eq:gamma_asymp}
  \gamma(t)
  \approx \frac{\Delta p_2}{\Delta p_1\,L_0}\,
    \frac{2\,\alpha_*\sqrt{\chi_0\,t}}{r_w} \to 0.
\end{equation}

\subsection{Начальная газонасыщенность}

Квазиравновесие (flash):
\begin{equation}\label{eq:Sgw_equil}
  \boxed{
    S_{g,w}^{\mathrm{eq}}
    = \frac{[R_s(p_b, T_0) - R_s(p_{wf}, T_0)]\,
      B_g(p_{wf}, T_0)}{B_o(p_{wf}, T_0)}.
  }
\end{equation}

$S_{g,w}^{\mathrm{eq}}$ устанавливается при $t = 0^+$
(локальное равновесие Black~Oil, $T = T_0$).
Формула~\eqref{eq:Sgw_equil} --- линеаризация при $S_g \ll 1$;
точное выражение содержит $B_o + \Delta R_s\,B_g$
в знаменателе вместо~$B_o$.
Конвективный перенос газа потоком нефти
при $S_g < S_{g,\mathrm{cr}}$ снижает $S_{g,w}$,
делая~\eqref{eq:Sgw_equil} верхней оценкой.

\subsection{Начальная температура}

При $t = t_\varepsilon$: $\Delta T_w(t_\varepsilon) \approx 0$.
Температура ещё не успела измениться.

Первое ненулевое приближение: охлаждение определяется
балансом эффекта Джоуля--Томсона и теплоты дегазации.
Интегрируя~\eqref{eq:energy} по малому объёму
$r_w < r < r_w + \delta$ при $T'_w = 0$:
\begin{equation}\label{eq:dTw_init}
  (\rho c_p)_\eff\,\frac{\dd T_w}{\dd t}\bigg|_{t=0^+}
  \approx Q_\mathrm{JT}(r_w) + Q_\mathrm{ad}(r_w) + Q_\mathrm{deg}(r_w).
\end{equation}
Оба источника конечны при $t = 0^+$ (давление и его
градиент определены из erfc-решения), поэтому
$\Delta T_w(t_\varepsilon) = O(t_\varepsilon) \approx 0$
при $t_\varepsilon \sim 10^{-6}\text{--}10^{-4}$~с.

\subsection{Сводка: начальные данные для ODE-системы}

Для запуска интегрирования системы~\eqref{eq:ode_system}
при $t = t_\varepsilon > 0$:
\begin{equation}\label{eq:ic_regularized}
  \boxed{
    r_b(t_\varepsilon) = r_w + 2\,\alpha_*\sqrt{\chi_0\,t_\varepsilon},
    \quad
    S_{g,w}(t_\varepsilon) = S_{g,w}^{\mathrm{eq}},
    \quad
    \Delta T_w(t_\varepsilon) = 0.
  }
\end{equation}

$t_\varepsilon$ выбирается из условия $\delta/r_w \sim 10^{-2}\text{--}10^{-1}$,
что при $\chi_0 \sim 0{,}5$~м$^2$/с, $\alpha_* \approx 0{,}5$,
$r_w = 0{,}1$~м даёт $t_\varepsilon \sim 10^{-6}\text{--}10^{-4}$~с.
На этих временах $\Delta T_w \approx 0$: тепловой отклик пренебрежимо мал.


%=============================================================================
\section{Обсуждение}\label{sec:discussion}
%=============================================================================

\subsection{Преимущества подхода}

\begin{enumerate}[nosep]
  \item Три PDE редуцированы к системе трёх ODE ---
        вычислительная стоимость на порядки ниже.
  \item Подвижная граница $r_b(t)$ --- естественный параметр.
  \item Температурное поле включено в систему: обратная связь
        $T \to R_s(p,T)$ учитывается автоматически на каждом
        временном шаге.
  \item Интегральные характеристики (дебит, объём газа,
        средняя температура) воспроизводятся по построению.
\end{enumerate}

\subsection{Сложность решения ODE-системы}

Формально система~\eqref{eq:ode_system} состоит из трёх ODE,
что на порядки дешевле исходных PDE. Однако реальная сложность
определяется не размерностью, а структурой правых частей:

\begin{enumerate}[nosep]
  \item \textbf{Вычисление коэффициентов.}
        На каждом временном шаге необходимо вычислить 9 элементов
        матрицы $\mathcal{A}_{ij}$ (из которых 2 нулевых)
        и 3~компоненты $\mathcal{F}_i$. Каждый элемент ---
        линейная комбинация интегралов $I_\Phi^{(n)}$,
        $I_\Psi^{(n)}$, $I_S^{(n)}$, которые при замораживании
        PVT-свойств вычисляются аналитически.
        Стоимость одного шага --- $O(1)$.

  \item \textbf{Конвективный интеграл в уравнении энергии.}
        Подставляя $u_\alpha = -\lambda_\alpha\,\pp p/\pp r$
        и~\eqref{eq:dT_dr}:
        \[
          \int_{r_w}^{R_\ext} r\,(\rho c_p u)_\mathrm{conv}^*\,
            \frac{\pp T}{\pp r}\,\dd r
          = -\int_{r_w}^{R_\ext}
            \frac{(\rho c_p \lambda)_\mathrm{conv}^*}{r}\,
            \frac{\pp p}{\pp r}\,
            \Psi'(\zeta)\,\frac{\Delta T_w}{L_0}\,\dd r,
        \]
        где $(\rho c_p \lambda)_\mathrm{conv}^*
        = \sum_\alpha \rho_\alpha c_{p,\alpha}\,\lambda_\alpha
          - L\,\rho_g^\st\,(\lambda_o/B_o)\,\pp R_s/\pp T$.
        При замораживании коэффициентов и подстановке
        $\pp p/\pp r$ из~\eqref{eq:dp_inner}--\eqref{eq:dp_outer}
        интеграл сводится к комбинации $I_\Phi^{(n)}$
        и $J^{(m)}$ через связь $\zeta(\eta)$.

  \item \textbf{Квадратичный член в уравнении давления.}
        Интеграл $\int r\,(\lambda_o/B_o)\,\nabla R_s\cdot\nabla p\,
        |\nabla p|\,\dd r$ содержит $(\pp p/\pp r)^2$ и
        $(\pp T/\pp r)\,(\pp p/\pp r)$. Второе слагаемое ---
        произведение производных \emph{разных} профилей,
        определённых в \emph{разных} координатах
        ($\eta$ для $p$, $\zeta$ для $T$), что не сводится
        к стандартным $I$-интегралам.

  \item \textbf{Жёсткость.}
        Система потенциально жёсткая: масштаб $\dot{r}_b$
        определяется гидродинамикой ($\chi_0 \sim 0{,}5$~м$^2$/с),
        а $\dot{\Delta T}_w$ --- температуропроводностью
        ($a \sim 10^{-7}$~м$^2$/с). Отношение масштабов
        $\chi_0/a \sim 10^6$--$10^7$, что типично
        для жёстких систем.
        Конвективный перенос ($\mathrm{Pe} \gg 1$) смягчает
        жёсткость: эффективная температуропроводность
        $a_\eff \sim uL \sim 10^{-4}$~м$^2$/с
        ($u \sim 10^{-3}$~м/с, $L \sim r_w$), что
        снижает отношение до $\chi_0/a_\eff \sim 10^3$--$10^4$.
        Система остаётся жёсткой; необходим неявный
        интегратор (BDF порядка 2--5).

  \item \textbf{Расщепление.}
        Структура~\eqref{eq:matrix_structure} позволяет
        решать подсистему $2\times 2$ на $(r_b, \Delta T_w)$
        и скалярное уравнение на $S_{g,w}$ последовательно,
        что упрощает обращение матрицы.
\end{enumerate}

\paragraph{Обезразмеривание.}
Элементы $\mathcal{A}_{11}$ масштабируются как $C_p$~[Па$^{-1}$],
а $\mathcal{A}_{33}$ --- как $(\rho c_p)_\eff^*$~[Дж/(м$^3$\,К)].
Без обезразмеривания обусловленность матрицы $\sim 10^{15}$.
Рекомендуется ввести характерные масштабы
$[p] = p_0 - p_{wf}$, $[T] = 10$~К, $[r] = r_w$, $[t] = r_w^2/\chi_0$
и перейти к безразмерным переменным $\tilde{r}_b$,
$\tilde{\Delta T}_w$, $\tilde{S}_{g,w}$.

\paragraph{Итого.}
При замораживании PVT-свойств система~\eqref{eq:ode_system}
допускает решение BDF-методом.
Основная вычислительная стоимость --- не в решении ODE,
а в вычислении PVT-свойств при обновлении средних $\bar{p}$,
$\bar{T}$, $\bar{S}_g$. Если PVT заданы корреляциями (Стендинг
и т.\,п.), стоимость одного шага $\sim 10^2$ операций.


\subsection{Ограничения и открытые вопросы}

\paragraph{Решённые вопросы.}
\begin{itemize}[nosep]
  \item Профиль температуры~\eqref{eq:Psi} удовлетворяет
        $\Psi'(1) = 0$, что обеспечивает $T'(r_w) = 0$
        (адиабатическое ГУ на забое).
  \item Зависимость $p_b(T)$ учтена через
        $p_b^* = p_b\bigl(T(r_b)\bigr)$
        в профиле давления и ограничении~\eqref{eq:gamma_constraint}.
  \item Двухзонное замораживание PVT-свойств
        (раздельные средние для зон разгазирования и однофазной)
        снижает ошибку по сравнению с единым средним.
  \item Смешанный интеграл $\int r\,(\pp T/\pp r)\,(\pp p/\pp r)\,\dd r$
        вычисляется аналитически благодаря линейной связи
        $\zeta = 1 - (L_b/L_0)\eta$.
  \item Модифицированные источники $Q_\mathrm{JT}^*$,
        $Q_\mathrm{ad}^*$ выписаны
        явно~\eqref{eq:QJT_star}--\eqref{eq:Qad_star}.
  \item Конвективный интеграл в уравнении энергии
        выражен через замороженные коэффициенты
        и стандартные $I$-интегралы.
\end{itemize}

\paragraph{Открытые вопросы.}
\begin{enumerate}[nosep]
  \item \textbf{Квадратичный член $(\pp p/\pp r)^2$.}
        Интеграл $\int r\,(\lambda_o/B_o)\,(\pp R_s/\pp p)\,
        (\pp p/\pp r)^2\,\dd r$ не сводится к граничным потокам
        и требует замораживания $\lambda_o/B_o$ при среднем давлении.
        Его вклад может быть оценён a~posteriori.

  \item \textbf{Перегруппировка $Q_\mathrm{deg}$.}
        Перегруппированные коэффициенты~\eqref{eq:rhocp_eff_mod},
        \eqref{eq:rhocp_conv_mod} зависят от $p$, $T$, $S_g$,
        что усложняет замораживание. При двухзонной оценке
        эффект смягчается, но строгая оценка погрешности открыта.

  \item \textbf{Выбор $\beta$, $\sigma$.}
        Показатели формы профилей можно зафиксировать или определить
        из старших моментов. Чувствительность подлежит исследованию.

  \item \textbf{Верификация.}
        Сопоставление с численным решением полной PDE-системы
        (метод конечных разностей или OPM Flow)
        необходимо для оценки реальной точности метода моментов.
\end{enumerate}


%=============================================================================
\section*{Литература}
%=============================================================================

\begin{thebibliography}{9}

\bibitem{draft4}
  Двухфазная фильтрация нефти с газом в радиальной геометрии.
  Полная система уравнений с балансом энергии.
  --- Draft4/two\_phase\_energy.tex (рабочий документ).

\bibitem{barenblatt1954}
  Баренблатт~Г.\,И. О некоторых приближённых методах в теории
  одномерной неустановившейся фильтрации жидкости при упругом
  режиме~// Изв. АН СССР, ОТН. --- 1954. --- №\,9. --- С.\,35--49.

\bibitem{BER1984}
  Баренблатт~Г.\,И., Ентов~В.\,М., Рыжик~В.\,М.
  Движение жидкостей и газов в природных пластах. ---
  М.: Недра, 1984. --- 211~с.

\bibitem{mirz2004}
  Мирзаджанзаде~А.\,Х., Хасанов~М.\,М., Бахтизин~Р.\,Н.
  Моделирование процессов нефтегазодобычи. ---
  М.--Ижевск: Институт компьютерных исследований, 2004. ---
  С.\,212--218 (§\,3.7).

\bibitem{pirverdyan1961}
  Пирвердян~А.\,М. Об одном способе оценок приближённых решений
  уравнений нестационарной фильтрации нефти и газа~//
  ПММ. --- 1961. --- Т.\,25, вып.\,4.

\bibitem{entov1965}
  Ентов~В.\,М. Теоремы сравнения для уравнений нестационарной
  фильтрации~// ПММ. --- 1965. --- Т.\,29, вып.\,11.

\end{thebibliography}

\end{document}
